%======================================================================
\chapter{PV-Knoten-Integration im Tensor Power Flow}
\label{ch:methoden_pv}
%======================================================================


Die zentrale Problemstellung verbunden mit der fehlenden Möglichkeit, PV-Knoten in dem Tensor Power Flow zu berücksichtigen wurde nun in \cref{ch:tpf} aufgezeigt. 
Dabei ist aufgewiesen, dass der TPF durch seine Tensorisierung der Leistungsgrößen einen erheblichen Vorteil gegenüber dem NR-Verfahren bezüglich der Rechenzeit haben kann.
Dieser wird darunter auch durch die sich nicht verändernde Systemmatrix $\tens{F}$ begründet.
Aufgrund dieser Eigenschaft muss dieser Tensor nur einmal zu Beginn der Vorschrift berechnet werden, vorausgesetzt das elektrische Netz besteht nur aus Slack- und PQ-Knoten.
Verfügt das Netz jedoch auch über PV-Knoten, beinhaltet die Iterationsvorschrift \eqref{eq:tpf_konstP} unbekannte Blindleistungen, welche das Verfahren nicht liefern kann.\\

In diesem Kapitel soll nun für diese unbekannte Größe eine weitere Iterationsvorschrift entwickelt werden, welche in Form einer äußeren Schleife mehrere TPF Durchgänge startet und die Blindleistung jedes mal neu schätzt.
Ziel dieser Methode soll es sein, den Effizienzvorteil des TPFs möglichst wenig zu verringern und die Tensorisierung des FPIs aufrecht zu erhalten.
Zunächst wird konkret analysiert, welche Auswirkungen die Inklusion von PV-Knoten auf die Iterationsgrößen haben.
Daraufhin wird unter diesen Informationen eine äußere Schleife konstruiert, welche auf Basis einer Thévenin-Näherung funktioniert.
Das Kapitel schließt mit einer Tensorformulierung der äußeren Schleife ab.

%======================================================================
\section{Analyse der Auswirkung von PV-Knoten auf die FPI-Struktur}
\label{sec:pv_auswirkung}
%======================================================================


Eine möglichst effiziente Implementierung der äußeren Schleife benötigt zunächst eine Analyse über die Auswirkung eines nicht-konstanten Leistungsvektors.
Dabei ist der Leistungsvektor

\begin{equation}
	s_k^{(n)} = P_k^{\spec} + \jj Q_k^{(n)}  
	\label{eq:s_pv_change}
\end{equation}

über die unbekannte Blindleistung $Q_\mathrm{k}^{(n)}$ der Iteration $n$ direkt betroffen. 
Deswegen müssen alle Größen der Iterationsvoschrift auf Ihre Abhängigkeit von der Blindleistung untersucht werden. 
Hierzu werden die Matrizen und Vektoren \eqref{eq:tpf_A}--\eqref{eq:tpf_w} in Tabelle \ref{tab:pv_abhaengigkeit_allgemein} mit ihrer Abhängigkeit von der Leistung aufgelistet.\\




%Um die nachfolgenden Methoden zu motivieren, wird zunächst systematisch analysiert, welche Bestandteile der FPI-Formulierung von einer Änderung der Blindleistung $Q_k$ an einem PV-Knoten $k$ betroffen sind. Der Leistungsvektor ändert sich am PV-Knoten gemäß:
%
%
%Die Auswirkungen auf die einzelnen Matrizen sind in \cref{tab:pv_abhaengigkeit_allgemein} zusammengefasst.

\begin{table}[H]
	\centering
	\caption{Abhängigkeit der FPI-Matrizen von der Blindleistung $Q_k$ eines PV-Knotens für den allgemeinen ZIP-Fall.}
	\label{tab:pv_abhaengigkeit_allgemein}
	\begin{tabular}{@{}lcc@{}}
		\toprule
		\textbf{Größe} & \textbf{Betroffenheit} & \textbf{Bedingung für Konstantheit}\\
		\midrule
		$\Ydd$ & Nein & Immer konstant  \\
		$\vect{A} = \diag(\alpha_P \had \vect{s}^*)$ & Ja & $\alpha_{P,k} = 0$ \\
		$\vect{B} = \diag(\alpha_Z \had \vect{s}^*) + \Ydd$ & Ja & $\alpha_{Z,k} = 0$ \\
		$\vect{c} = \Yds\vect{v}_s + \alpha_I \had \vect{s}^*$ & Ja & $\alpha_{I,k} = 0$ \\
		$\vect{B}^{-1}$ & Ja & $\alpha_{Z,k} = 0$  \\
		$\vect{F} = -\vect{B}^{-1}\vect{A}$ & Ja & $\alpha_{P,k} = \alpha_{Z,k} = 0$ \\
		$\vect{w} = -\vect{B}^{-1}\vect{c}$ & Ja & $\alpha_{Z,k} = \alpha_{I,k} = 0$  \\
		\bottomrule
	\end{tabular}
\end{table}

Dabei ist zu erkennen, dass der ZIP-Fall der tragende Faktor für die Konstanz der Matrizen ist.
Ist ein Vollständiger ZIP-Fall gegeben ist weder $\mathbf{F}$ noch $\mathbf{w}$ konstant über alle Iterationen.
Da das Interesse bei der Systemmatrix liegt, lassen sich hier zwei Fälle bestimmen.\\

Wenn $\alpha_{Z,k} = 0$ ist, bleiben die für $\mathbf{F}$ relevante Matrix $\mathbf{B}$ konstant und der einzige Leistungseinfluss bleibt mit $\alpha_{P,k}$, dem konstanten Lastanteil. Dieser Fall hat als Vorteil, dass die $\mathbf{B}$-Matrix nur einmal invertiert werden muss.
Ist $\alpha_{Z,k} \neq 0$ ist die Konstanz nicht mehr gegeben und der Rechenaufwand erhöht sich durch die sich wiederholende Inversion von $\mathbf{B}$.
Im konstanten Stromfall $\alpha_{I,k} \neq 0$ bleibt der Rechenvorteil der einmaligen Inversion erhalten. Ferner verschwindet in diesem Fall die Systemmatrix $\mathbf{F}$ und die einzige Veränderliche ist der Slack-Einfluss $\mathbf{w}$.\\

Generatoren werden häufig, wie in \cref{ch:grundlagen} beschrieben als konstante Leistungsspeiser implementiert, was den konstanten Lastfall zum Fokus und Standardfall dieser Arbeit macht. 
Weiterhin folgt dieser Analyse, dass die äußere Blindleistungskorrektur zunächst ausschließlich PV-Zeilen umfassen muss.
Dies ist eine direkte Konsequenz der $\operatorname{diag}$-Operation in der Systemmatrix, welche die Leistung knotenweise in die Rechnung einbringt.
Da die einzige Kopplung aus der der Bus-Impedanzmatrix folgt, kann eine Thévenin-Näherung verwendet werden, da alle relevanten Größen hierfür bereits bekannt sind und sich nur mit der Netztopologie Ändern würden. 
Damit lässt sich die Alternative einer arbeitspunktabhängigen Korrektur vermeiden, welche die Nachteile einer NR-Iteration mit sich bringen würde (siehe \cref{sec:nr}).
Abschließend soll die Korrektur selber tensorisierbar über alle Lastflüsse $\tau$ sein, um der in dem TPF zugrunde gelegte Parallelisierung auszunutzen.\\


%
%
%
%%\paragraph{Interpretation.}
%Die Auswirkungen einer Q-Änderung am PV-Knoten hängen fundamental vom Lastmodell an diesem Knoten ab:
%\begin{itemize}
%	\item \textbf{Fall 1:} $\alpha_{Z,k} = 0$ am PV-Knoten. Dann bleibt $\vect{B}$ (und damit $\vect{B}^{-1}$) über alle Iterationen \textit{konstant}. Lediglich $\vect{A}$ und ggf.\ $\vect{c}$ ändern sich. Dies ist der rechentechnisch günstigste Fall, da die Inversion von $\vect{B}$ nur einmal durchgeführt werden muss.
%	\item \textbf{Fall 2:} $\alpha_{Z,k} \neq 0$ am PV-Knoten. Dann ändert sich $\vect{B}$ mit jeder Q-Aktualisierung, und $\vect{B}^{-1}$ muss grundsätzlich neu berechnet werden. %Jedoch ändert sich nur ein \textit{Diagonalelement} von $\vect{B}$, was effiziente Rang-1-Updates (Sherman-Morrison-Formel) ermöglicht.
%\end{itemize}
%
%In der Praxis werden Generatoren an PV-Knoten fast ausschließlich als leistungskonstante Einspeiser modelliert ($\alpha_{P,k} = 1$, $\alpha_{Z,k} = \alpha_{I,k} = 0$), da ihr Verhalten durch die Leistungselektronik oder Turbinenregelung bestimmt wird, nicht durch die Netzspannung \cite{Garcia2000, Giraldo2022}. Für diesen häufigen Spezialfall gilt:
%\begin{equation}
%	\alpha_{Z,k} = \alpha_{I,k} = 0 \quad \Rightarrow \quad \vect{B}^{-1} \text{ und } \vect{w} \text{ bleiben konstant.}
%	\label{eq:pv_spezialfall}
%\end{equation}
%
%Die nachfolgende Entwicklung wird zunächst für den \textit{allgemeinen ZIP-Fall} formuliert und die Vereinfachungen für den Spezialfall \eqref{eq:pv_spezialfall} jeweils explizit angegeben.


%\subsection{Konsequenzen für den Konvergenzbeweis}
%\label{subsec:pv_konvergenz_konsequenz}

%Der Konvergenzbeweis nach Banach basiert auf der Konstantheit von $\vect{s}$ und damit aller Hilfsmatrizen. Wenn $\vect{s}$ sich zwischen den Iterationen ändert, ist der Kontraktionsfaktor $\eta$ aus \eqref{eq:eta} nicht mehr über die gesamte Iteration konstant.
%%
%%	\item Der formale Konvergenzbeweis (\cref{prop:sam_kontraktion}) ist in seiner vorliegenden Form \textit{nicht direkt anwendbar}, da die Abbildung $T^{(n)}(\cdot)$ iterationsabhängig wird.
%%	\item Die physikalische Konvergenzbedingung $\|Z_{ii}\| < \|Z_{L,i}\|$ (\cref{subsec:konvergenz_physik}) bleibt als heuristische Richtlinie anwendbar, sofern die Q-Änderungen die äquivalente Lastimpedanz nicht fundamental verändern.
%In der Praxis konvergieren die nachfolgend vorgestellten Methoden für typische Verteilnetz-Betriebspunkte zuverlässig, wie in \cref{ch:ergebnisse} numerisch validiert wird.


%======================================================================
\section{Methode der Äußeren Q-Schleife}
\label{sec:methode_a}
%======================================================================

%\subsection{Grundidee}
%\label{subsec:methode_a_idee}



Nachdem die strukturellen Anforderungen an die Methode spezifiziert wurden, gilt es den Aufbau der Schleife zu konkretisieren.
Das Grundkonzept ist hierbei eine Verschachtelung des TPFs mit der Q-Korrektur.
Der TPF wird ab hier als \emph{innere} Schleife bezeichnet und durchgeführt, während die Q-Korrektur als \emph{äußere} Schleife  den TPF ummantelt.
Während einer äußeren Iteration $l$ wird mit einem konstanten Leistungsvektor $\mathbf{s}^{(l)}$ gerechnet, bis das Verfahren wie in \cref{ch:tpf} spezifiziert konvergiert.
Nachfolgend wird anhand dieses Ergebnisses wird ein Blindleistungsvektor $\mathbf{q}$ für ausschließlich PV-Knoten korrigiert.
Das Ergebnis wird abschließend für den Leistungsvektor $\mathbf{s}{(\ell+1)}$ eingesetzt\\


%Methode~A realisiert die PV-Knoten-Behandlung durch eine verschachtelte
%Schleifenstruktur. Der Standard-TPF wird als \textit{innere Schleife} mit
%festem Leistungsvektor $\vect{s}^{(\ell)}$ (und damit festen Matrizen
%$\vect{F}^{(\ell)}$, $\vect{w}^{(\ell)}$) ausgeführt, bis er konvergiert.
%Anschließend wird in einer \textit{äußeren Schleife} der Blindleistungs\-
%vektor $\vect{q}\in\mathbb{R}^{n_{pv}}$ an \emph{allen} PV-Knoten simultan
%korrigiert, der Leistungsvektor $\vect{s}^{(\ell+1)}$ aktualisiert und die
%Hilfsmatrizen neu berechnet.

Der Ablauf soll nach dem folgenden Schema
\begin{enumerate}
	\item Setze $\vect{q}^{(0)} = \vect{0}$ als initialen Schätzwert
	(äußere Iteration $\ell = 0$).
	\item Berechne $\vect{s}^{(\ell)}$, $\vect{A}^{(\ell)}$,
	$\vect{B}^{(\ell)}$, $\vect{c}^{(\ell)}$, $\vect{F}^{(\ell)}$,
	$\vect{w}^{(\ell)}$.
	\item Führe den vollständigen TPF (innere Schleife) mit festem
	$\vect{F}^{(\ell)}$, $\vect{w}^{(\ell)}$ bis zur Konvergenz aus,
	welches zu $\vect{v}^{(\ell)}$ führt.
	\item Berechne den Vektor der Spannungsabweichungen
	\[
	\Delta|\vect{V}_{pv}|^2
	\;=\; |\vect{V}^{\spec}|^2 - |\vect{v}^{(\ell)}_{pv}|^2.
	\]
	\item Korrigiere $\vect{q}^{(\ell+1)} = \vect{q}^{(\ell)} + \Delta\vect{q}^{(\ell)}$
	{gekoppelt} über alle PV-Knoten
	(\cref{subsec:methode_a_herleitung}).
	%\item Falls $\bigl\lVert\Delta|\vect{V}_{pv}|^2\bigr\rVert_\infty < \varepsilon_Q$:
	%fertig. Sonst: zurück zu Schritt~2.
\end{enumerate}

ablaufen.
Bei diesem Ablauf ist der Vorteil, dass die innere Schleife einen unveränderten TPF beinhaltet. Die Matrizen und Vektoren im TPF bleiben über die innere Schleife konstant. Somit bleibt die robuste Konvergenz nach Banach, wie in \cref{subsec:banach} eingeführt, bestehen.
Damit werden alle weitere Effekte auf die äußere Schleife isoliert.
Nachfolgend soll die Sensitivität für die Blindleistungskorrektur formuliert werden, wobei das Hauptziel eine möglichst aufwandsarme Berechnung ist.



%Ein wesentlicher Vorteil dieser Struktur ist, dass die innere Schleife ein
%\textit{unmodifizierter} Standard-TPF ist. Da $\vect{s}^{(\ell)}$ dort
%konstant gehalten wird, sind alle Hilfsmatrizen über die inneren Iterationen
%fixiert. Der Konvergenzbeweis (\cref{eq:convergence_condition}) ist daher für
%die innere Schleife vollständig gültig.

%\subsection{Herleitung der gekoppelten Q-Korrekturformel}
%\label{subsec:methode_a_herleitung}


%
%
%
%
%
%Für die Q-Korrektur in Schritt~5 wird eine Beziehung zwischen dem
%\emph{Vektor} $\Delta\vect{q}$ und dem Vektor der Spannungsabweichungen
%$\Delta|\vect{V}_{pv}|^2$ benötigt. Wir leiten diese über das
%\emph{multivariate} Thévenin-Äquivalent am gesamten PV-Bus-Vektor her --
%d.\,h.\ unter voller Berücksichtigung der PV--PV-Kopplung über gemeinsame
%Netzimpedanzen.
%
%%\paragraph{Notation.}
%%Sei $\mathcal{P} = \{k_1, \ldots, k_{n_{pv}}\} \subseteq \{1, \ldots, b\varphi\}$
%%die Indexmenge der PV-Knoten im $d$-Block. Wir definieren
%%\begin{align*}
%%	\vect{V}_{pv} &= \vect{V}[\mathcal{P}] \in \mathbb{C}^{n_{pv}}
%%	&&\text{(Spannungen an PV-Knoten)} \\
%%	\vect{I}_{pv} &= \vect{I}[\mathcal{P}] \in \mathbb{C}^{n_{pv}}
%%	&&\text{(Strominjektionen)} \\
%%	\vect{q}     &\in \mathbb{R}^{n_{pv}}
%%	&&\text{(Blindleistungen der PV-Knoten)} \\
%%	\vect{Z}_B^{(\ell)}   &= \bigl(\vect{B}^{(\ell)}\bigr)^{-1}
%%	\in \mathbb{C}^{b\varphi\times b\varphi}
%%	&&\text{(Bus-Impedanzmatrix)} \\
%%	\vect{Z}_{pp}^{(\ell)} &= \vect{Z}_B^{(\ell)}[\mathcal{P}, \mathcal{P}]
%%	\in \mathbb{C}^{n_{pv}\times n_{pv}}
%%	&&\text{(PV--PV-Submatrix)} \\
%%	\vect{X}_{pp}^{(\ell)} &= \im\!\left(\vect{Z}_{pp}^{(\ell)}\right)
%%	\in \mathbb{R}^{n_{pv}\times n_{pv}}
%%	&&\text{(PV--PV-Reaktanz-Block)}
%%\end{align*}
%
%\paragraph{Multivariates Thévenin-Äquivalent.}
%Ausgehend von der Netzgleichung
%$\vect{B}^{(\ell)}\vect{V} + \vect{Y}_{ds}\vect{v}_s = \vect{I}$ folgt durch
%Umstellen
%\begin{equation}
%	\vect{V} \;=\; \vect{V}^{\text{open}} \;+\; \vect{Z}_B^{(\ell)}\,\vect{I},
%	\qquad
%	\vect{V}^{\text{open}} = -\vect{Z}_B^{(\ell)}\,\vect{Y}_{ds}\,\vect{v}_s,
%	\qquad
%	I_k = \frac{s_k^*}{V_k^*}
%	\label{eq:multivar_thevenin}
%\end{equation}
%mit $\vect{V}^{\text{open}}$ als Leerlaufspannungsvektor. Einschränkung auf
%die PV-Zeilen liefert das gekoppelte Thévenin-System:
%\begin{equation}
%	\vect{V}_{pv} \;=\; \vect{V}^{\text{open}}_{pv}
%	\;+\; \vect{Z}_{pp}^{(\ell)}\,\vect{I}_{pv}
%	\;+\; \sum_{j \notin \mathcal{P}} \vect{Z}_B^{(\ell)}[\mathcal{P}, j]\,I_j
%	\label{eq:pv_thevenin_block}
%\end{equation}
%Der letzte Term hängt nur von den PQ-Strömen ab und bleibt in erster
%Näherung während einer Q-Variation an den PV-Knoten konstant.
%
%\medskip
%%Die entscheidende Beobachtung im Vergleich zur ursprünglichen Herleitung ist:
%%\begin{center}
%%	\emph{$\vect{Z}_{pp}^{(\ell)}$ ist im Allgemeinen \emph{keine} Diagonalmatrix.}
%%\end{center}
%Ein Off-Diagonalelement $\vect{Z}_{pp}^{(\ell)}[k, j]$ mit $k \neq j$
%beschreibt, wie sich der Strom an PV-Knoten~$j$ auf die Spannung an
%PV-Knoten~$k$ auswirkt. In radialen Netzen entspricht
%$\vect{Z}_{pp}^{(\ell)}[k, j]$ der \emph{Summe der Serienimpedanzen entlang
%	des gemeinsamen Vorfahrenpfads} zwischen den Knoten~$k$ und~$j$ (ausgehend
%vom Slack). Zwei PV-Knoten in demselben Ast des Baumes haben daher eine
%besonders starke Kopplung.
%
%\paragraph{Linearisierung: Stromänderung bei Q-Variation.}
%Bei festgehaltener Wirkleistung $P_j$ gilt am PV-Knoten~$j$
%$I_j = (P_j + \jj q_j)^\ast/V_j^*$ (Erzeugerkonvention). Linearisierung um
%den aktuellen Arbeitspunkt $v_j^{(\ell)}$ ergibt
%\begin{equation}
%	\Delta I_j = \frac{-\jj\,\Delta q_j}{V_j^*}\bigg|_{V_j = v_j^{(\ell)}} .
%	\label{eq:delta_I_j}
%\end{equation}
%In kompakter Vektorform:
%\begin{equation}
%	\Delta\vect{I}_{pv} = -\jj\,\vect{D}_{V^*}^{-1}\,\Delta\vect{q},
%	\qquad
%	\vect{D}_{V^*} := \diag\!\left(v_{k_1}^{(\ell)\,*}, \ldots,
%	v_{k_{n_{pv}}}^{(\ell)\,*}\right) .
%	\label{eq:delta_I_vec}
%\end{equation}
%
%\paragraph{Gekoppelte Spannungsänderung.}
%Da alle Terme in~\eqref{eq:pv_thevenin_block} außer
%$\vect{Z}_{pp}^{(\ell)}\vect{I}_{pv}$ in der Linearisierung konstant sind:
%\begin{equation}
%	\Delta\vect{V}_{pv} \;=\; \vect{Z}_{pp}^{(\ell)}\,\Delta\vect{I}_{pv}
%	\;=\; -\jj\,\vect{Z}_{pp}^{(\ell)}\,\vect{D}_{V^*}^{-1}\,\Delta\vect{q}
%	\label{eq:delta_V_coupled}
%\end{equation}
%
%Komponentenweise für $k \in \mathcal{P}$:
%\begin{equation}
%	\Delta V_k
%	\;=\; \sum_{j \in \mathcal{P}}
%	\underbrace{\vect{Z}_{pp}^{(\ell)}[k, j]}_{\text{gekoppelt}}
%	\cdot \frac{-\jj\,\Delta q_j}{v_j^{(\ell)\,*}}
%	\label{eq:delta_V_component}
%\end{equation}
%Im Gegensatz zur klassischen (entkoppelten) Herleitung wirken \emph{alle}
%$\Delta q_j$ simultan auf $\Delta V_k$: jede Q-Änderung an einem PV-Knoten
%koppelt in jede andere PV-Knoten-Spannung ein.
%
%\paragraph{Änderung des Spannungsbetragsquadrats.}
%Mit der Identität
%$\Delta(|V_k|^2) = 2\,\re\!\left(V_k^*\,\Delta V_k\right)$
%und $V_k^* = |v_k^{(\ell)}|\,e^{-\jj\theta_k}$,
%$1/V_j^* = e^{\jj\theta_j}/|v_j^{(\ell)}|$ folgt aus~\eqref{eq:delta_V_component}
%\begin{align}
%	V_k^*\,\Delta V_k
%	&= |v_k^{(\ell)}|\,e^{-\jj\theta_k}\cdot
%	\sum_{j \in \mathcal{P}} \vect{Z}_{pp}^{(\ell)}[k, j]\cdot
%	\frac{-\jj\,e^{\jj\theta_j}}{|v_j^{(\ell)}|}\,\Delta q_j \nonumber \\
%	&= -\jj\,\sum_{j \in \mathcal{P}} \vect{Z}_{pp}^{(\ell)}[k, j]\cdot
%	\underbrace{\frac{|v_k^{(\ell)}|}{|v_j^{(\ell)}|}\,
%		e^{\jj(\theta_j - \theta_k)}}_{\approx\,1}\cdot \Delta q_j .
%	\label{eq:VstarDeltaV_coupled}
%\end{align}
%Unter der Annahme $|v_k^{(\ell)}| \approx |v_j^{(\ell)}|$ und kleiner
%Winkeldifferenzen $\theta_j - \theta_k \approx 0$ (typisch für radiale
%Verteilnetze im Normalbetrieb) reduziert sich der Faktor zu~$1$. Mit
%$\vect{Z}_{pp}^{(\ell)}[k, j] = R_{kj}^{(\ell)} + \jj X_{kj}^{(\ell)}$ folgt
%\begin{equation}
%	V_k^*\,\Delta V_k
%	\;\approx\; \sum_{j \in \mathcal{P}}
%	\bigl(X_{kj}^{(\ell)} - \jj R_{kj}^{(\ell)}\bigr)\,\Delta q_j
%	\label{eq:VstarDeltaV_simplified}
%\end{equation}
%und damit für den Realteil
%\begin{equation}
%	\re\!\left(V_k^*\,\Delta V_k\right)
%	\;=\; \sum_{j \in \mathcal{P}} X_{kj}^{(\ell)}\,\Delta q_j
%	\;=\; \bigl[\vect{X}_{pp}^{(\ell)}\,\Delta\vect{q}\bigr]_k .
%	\label{eq:real_part_coupled}
%\end{equation}
%Bemerkenswert: Wie im skalaren Fall entfallen \emph{sämtliche}
%Widerstandsanteile $R_{kj}^{(\ell)}$ vollständig -- auch die
%Off-Diagonalen mit $k \neq j$.
%
%\paragraph{Zentrale Sensitivitätsbeziehung (Matrixform).}
%Zusammenfassend gilt für den gesamten PV-Vektor:
%\begin{equation}
%	\boxed{\;\Delta|\vect{V}_{pv}|^2 \;=\; 2\,\vect{X}_{pp}^{(\ell)}\,\Delta\vect{q}\;}
%	\label{eq:sensitivity_matrix}
%\end{equation}
%mit $\vect{X}_{pp}^{(\ell)} \in \mathbb{R}^{n_{pv}\times n_{pv}}$.
%Für $n_{pv} = 1$ ist $\vect{X}_{pp}^{(\ell)}$ eine $1\times 1$-Matrix mit
%dem einzigen Eintrag $X_{kk}^{(\ell)}$, und~\eqref{eq:sensitivity_matrix}
%reduziert sich exakt auf $\Delta(|V_k|^2) = 2 X_{kk}^{(\ell)}\,\Delta q_k$.
%Die klassische entkoppelte Formel ist somit ein Sonderfall.
%
%\paragraph{Ergebnis: gekoppelte Q-Korrekturformel.}
%Auflösen nach $\Delta\vect{q}$ mit Zielwert
%$\Delta|\vect{V}_{pv}|^2 = |\vect{V}^{\spec}|^2 - |\vect{v}^{(\ell)}_{pv}|^2$
%liefert die gekoppelte Q-Korrektur
%\begin{equation}
%	\boxed{\;
%		\Delta\vect{q}^{(\ell)}
%		\;=\; \tfrac{1}{2}\,\bigl(\vect{X}_{pp}^{(\ell)}\bigr)^{-1}
%		\bigl(|\vect{V}^{\spec}|^2 - |\vect{v}^{(\ell)}_{pv}|^2\bigr) .
%		\;}
%	\label{eq:methode_a_dQ_coupled}
%\end{equation}
%Dies ersetzt die entkoppelte Version
%\begin{equation}
%	\Delta q_k^{(\ell)}
%	\;=\; \frac{|V_k^{\spec}|^2 - |v_k^{(\ell)}|^2}{2\,X_{kk}^{(\ell)}}
%	\qquad\text{(entkoppelt, nur Diagonale)}
%	\label{eq:methode_a_dQ_scalar}
%\end{equation}
%durch eine simultane Lösung des \emph{gesamten} PV-Systems: jedes
%$\Delta q_j$ berücksichtigt die Rückwirkung auf alle anderen PV-Knoten via
%$\vect{X}_{pp}^{(\ell)}$.
%
%\paragraph{Numerische Umsetzung.}
%Sind alle $\alpha_{Z, k} = 0$ an den PV-Knoten, so ist
%$\vect{B}^{(\ell)} = \vect{B}$ konstant über die äußere Iteration. Damit
%lässt sich $\vect{X}_{pp}$ (bzw.\ dessen LU-Faktorisierung oder Inverse)
%\emph{einmalig} vor Beginn der äußeren Schleife berechnen. Die Auswertung
%von~\eqref{eq:methode_a_dQ_coupled} in jeder Iteration reduziert sich auf
%ein Matrix-Vektor-Produkt der Kosten $\mathcal{O}(n_{pv}^2)$, was
%gegenüber den inneren FPI-Kosten
%$\mathcal{O}\bigl((b\varphi)^2\,\tau\bigr)$ vernachlässigbar bleibt
%($n_{pv} \ll b\varphi$). Die Tensorstruktur wird nicht aufgebrochen:
%$(\vect{X}_{pp})^{-1}$ wird über alle $\tau$ Szenarien broadcast.



\subsection{Herleitung der Q-Korrekturformel}
\label{subsec:methode_a_herleitung}


In Schritt~5 wird eine Spannungsdifferenz zwischen dem durch PV-Knoten gegebenen Spannungsbetrag und dem tatsächlich ausgerechneten Spannungsabfall berechnet.
Diese Differenz soll als Residuum zur Blindleistungskorrektur dienen, da sich bei Angleichung der PV-Spannungen die nötige Blindleistung einstellt um die PV-Vorgabe zu halten.
Da das Residuum über das Kontraktionsverhalten zudem auch über die Konvergenzbedingung verfügt, wird nur noch ein Zusammenhang benötigt, welcher die Korrekturrichtung vorgibt, in welche die Q-Korrektur erfolgen soll.\\


Die Thévenin-Impedanz ist eine einfache Möglichkeit Spannung und Strom am Arbeitspunkt eines PV-Knoten in Beziehung zu setzen.
Diese kann folglich dafür verwendet werden, die Reaktion der Leistung auf auf eine Spannungsänderung zu approximieren.\\

%Für die Q-Korrektur in Schritt~5 wird eine Beziehung zwischen $\Delta Q_k$ und
%dem Spannungsmismatch $\Delta|V_k|^2$ benötigt. Diese wird über das
%Thévenin-Ersatzschaltbild am PV-Knoten~$k$ hergeleitet.

%----------------------------------------------------------------------
%\paragraph{Thévenin-Äquivalent am Knoten $k$.}
Gemäß \cref{subsec:thevenin} ist die das Thévenin-Äquivalent an einem PV-Knoten $k$ durch
\begin{equation}
	V_k \;=\; E_{th,k} \;-\; Z^{(\ell)}_{B,kk}\cdot I_k,
	\label{eq:thevenin_skalar}
\end{equation}
gegeben. Dabei ist $E_{th,k}$ die Leerlaufspannung des äquivalenten Schaltbildes und $Z^{(\ell)}_{B,kk}$ die Thévenin-Impedanz, wie es auch in Abbildung \ref{fig:thevenin} gezeigt wird. 
Die Impedanz ist darüber hinaus das jeweilige Diagonalelement der Impedanzmatrix 


%wobei $E_{th,k}$ die Thévenin-Leerlaufspannung (bestimmt durch das restliche Netz
%und die Slack-Spannung) und $Z^{(\ell)}_{B,kk}$ das $k$-te Diagonalelement der
%Bus-Impedanzmatrix ist:
\begin{equation}
	Z^{(\ell)}_{B,kk} \;=\; \bigl[\ZB^{(\ell)}\bigr]_{kk}
	\;=\; R^{(\ell)}_{kk} + \jj\,X^{(\ell)}_{kk}.
	\label{eq:diagonalelement}
\end{equation}

Für eine konstante Netztopologie, wie sie für die Arbeit angenommen wird, sind die Impedanzen über alle äußeren Iterationen $\ell$ gleich.
%----------------------------------------------------------------------
%\paragraph{Linearisierung: Stromänderung bei Q-Variation.}
%Bei festgehaltenem $P_k$ und unter der Annahme, dass sich $V_k$ bei einer kleinen
%Änderung $\Delta Q_k$ nur geringfügig ändert (Linearisierung um den aktuellen
%Arbeitspunkt $v^{(\ell)}_k$), ergibt sich die Stromänderung:
Von Interesse ist an der Stelle nicht der absolute Wert von der Spannung, sondern dessen Differenz.
Wird diese aus \eqref{eq:thevenin_skalar} gebildet, folgt der Ausdruck
\begin{equation}
	\Delta V_k \;=\; -Z^{(\ell)}_{B,kk}\cdot\Delta I_k
	\label{eq:spannung_skalar}
\end{equation}

mit noch einem unbekannten Ausdruck für die Stromdifferenz.
Der Knotenstrom 

%----------------------------------------------------------------------
%\paragraph{Strominjektion als Funktion von $Q_k$.}
%Die komplexe Leistungsinjektion am Knoten~$k$ in Erzeugerkonvention lautet

\begin{equation}
	I_k \;=\; \frac{S_k^*}{V_k^*} \;=\; \frac{P_k - \jj Q_k}{V_k^*}.
	\label{eq:strom_skalar}
\end{equation}
kann als Funktion der Knotenleistung $S_{k}$ dargestellt werden. Der  Ausduck folgt direkt aus der Leistungsdefinition \eqref{eq:power_current}.
Bei Bildung von der Stromänderung wird beachtet, dass die Wirkleistung $P_{k}$ gegeben ist und sich folglich eliminiert. In der Nähe des Arbeitspunktes $v_{k}$ folgt die Stromänderung aufgrund einer Blindleistungsinjektion

\begin{equation}
	\Delta I_k \;\approx\; \left.\frac{-\jj\,\Delta Q_k}{V_k^*}\right|_{V_k = v^{(\ell)}_k}.
	\label{eq:strom_linear}
\end{equation}

Dieser linearisierte Term kann weiterhin in Gleichung \eqref{eq:spannung_skalar} eingesetzt und sortiert werden. 
Die Spannungsänderung 

%----------------------------------------------------------------------
%\paragraph{Spannungsänderung über Thévenin-Impedanz.}
%Die resultierende Spannungsänderung am Knoten~$k$ folgt aus dem
%Thévenin-Äquivalent~\eqref{eq:thevenin_skalar}:
\begin{equation}
	V_k^* \cdot \Delta V_k \;=\; Z^{(\ell)}_{B,kk}\cdot\jj\,\Delta Q_k
	\label{eq:spannung_skalar}
\end{equation}

ist nach einsetzten somit eine Funktion des Arbeitspunkt und der Blindleistungsänderung.
Damit liegt bereits ein Ausdruck vor, mit welcher eine Blindleistungskorrektur ausgehend von einer Spannungsdifferenz in Beziehung stehen.
Der Nachteil dieser Formulierung ist jedoch dass die Spannungsdifferenz selber noch eine komplex Größe ist. 
Dies setzt folglich für eine Differenz zwischen aktueller und vorgegebener PV-Spannung voraus, dass die Winkeldifferenz bekannt ist.
Da aber eine Winkelvorgabe bei PV-Knoten fehlt, muss die Formulierung auf Spannungsbeträge erweitert werden, um die Winkelabhängigkeit zu verlieren.
Um das zu erzielen, kann das Produkt 	$V_k^* V_k = |V_k|^2$ gebildet werden.
Wendet man die Differenz auf das Betragsquadrat an an, kann der Ausdruck 

\begin{equation}
    \Delta(|V_k|^2) = |V_k + \Delta V_k|^2 - |V_k|^2
	= V_k^*\Delta V_k + V_k(\Delta V_k)^* + |\Delta V_k|^2
	\label{eq:V2}
\end{equation}

formuliert werden.
Da ohnehin die Linearisierung um $V_k$ betrachtet wird, kann $|\Delta V_k|^2$ vernachlässigt werden. Durch die Identität $z + z^* = 2\,\re(z)$ reduziert sich der Ausdruck \eqref{eq:V2} zu

\begin{equation}
	\Delta(|V_k|^2) \;=\; 2\,\re\bigl(V_k^*\cdot\Delta V_k\bigr)
	\label{eq:betragsquadrat_identitaet}
\end{equation}

Das Argument der Realteils entspricht genau der rechten Seite von Gleichung \eqref{eq:spannung_skalar}.
Dementsprechend kann dessen Realteil

\begin{equation}
	\re\bigl(V_k^*\cdot\Delta V_k\bigr) \;=\; X^{(\ell)}_{kk}\cdot\Delta Q_k.
	\label{eq:realteil_skalar}
\end{equation}

ausgewertet werden. 
Dabei ist die Reaktanz $X^{(\ell)}$ der Imaginärteil der Impedanzmatrix $Z^{(\ell)}_{B,kk}$ und trägt als einziger zum Realteil des Ausdrucks $Z^{(\ell)}_{B,kk}\cdot \jj$.
Zusammen mit \eqref{eq:betragsquadrat_identitaet} kann die Sensitivitätsbeziehung 

\begin{equation}
	\Delta(|V_k|^2) \;=\; 2\,X^{(\ell)}_{kk}\cdot\Delta Q_k.
	\label{eq:sensitivity_skalar}
\end{equation}

formuliert werden.
Dabei ist die Spannungsdifferenz gegeben durch den aktuell berechneten Spannungsbetrag $|V_k|$ und dem vorgegebenen PV-Wert $|V^{\spec}_k|$.
Eingesetzt in \eqref{eq:sensitivity_skalar} ist die Berechnungsvorschrift für die Blindleistungskorrektur durch

%Hierzu werden die Spannungsgrößen auf eine Seite 	$V_k^* V_k = $
%
%\begin{equation}
%	V_k^* \Delta V_k \;=\; Z^{(\ell)}_{B,kk}\cdot \jj\,\Delta Q_k
%	\label{eq:spannung_skalar}
%\end{equation}
%
%gebracht.
%
%
%
%Mit $Z^{(\ell)}_{B,kk} = R^{(\ell)}_{kk} + \jj X^{(\ell)}_{kk}$ und
%$(v^{(\ell)}_k)^* = |v^{(\ell)}_k|\,e^{-\jj\theta_k}$ ergibt die Ausmultiplikation:
%\begin{equation}
%	\Delta V_k \;=\; \frac{\Delta Q_k}{|v^{(\ell)}_k|}
%	\bigl(X^{(\ell)}_{kk} - \jj R^{(\ell)}_{kk}\bigr)\,e^{\jj\theta_k}.
%	\label{eq:spannung_ausmult}
%\end{equation}
%
%%----------------------------------------------------------------------
%%\paragraph{Änderung des Spannungsbetragsquadrats.}
%Die Änderung von $|V_k|^2$ lässt sich über die Identität\footnote{Diese folgt
%	direkt aus $|V_k|^2 = V_k V_k^*$: Setzt man den geänderten Wert
%	$V_k + \Delta V_k$ ein und multipliziert aus, erhält man $\Delta(|V_k|^2)
%	= V_k^*\Delta V_k + V_k(\Delta V_k)^* + |\Delta V_k|^2$. Der quadratische Term
%	$|\Delta V_k|^2$ wird im Sinne der Linearisierung vernachlässigt; mit
%	$z + z^* = 2\,\re(z)$ verbleibt $\Delta(|V_k|^2) = 2\,\re(V_k^*\Delta V_k)$.}
%\begin{equation}
%	\Delta(|V_k|^2) \;=\; 2\,\re\bigl(V_k^*\cdot\Delta V_k\bigr)
%	\label{eq:betragsquadrat_identitaet}
%\end{equation}
%berechnen. Einsetzen von~\eqref{eq:spannung_ausmult} mit
%$V_k^* = |v^{(\ell)}_k|\,e^{-\jj\theta_k}$ liefert:
%\begin{align}
%	V_k^*\cdot\Delta V_k
%	&= |v^{(\ell)}_k|\,e^{-\jj\theta_k}\cdot\frac{\Delta Q_k}{|v^{(\ell)}_k|}
%	\bigl(X^{(\ell)}_{kk} - \jj R^{(\ell)}_{kk}\bigr)\,e^{\jj\theta_k} \notag\\
%	&= \Delta Q_k\cdot\bigl(X^{(\ell)}_{kk} - \jj R^{(\ell)}_{kk}\bigr).
%	\label{eq:VstarDeltaV_skalar}
%\end{align}
%Der Realteil ergibt:
%\begin{equation}
%	\re\bigl(V_k^*\cdot\Delta V_k\bigr) \;=\; X^{(\ell)}_{kk}\cdot\Delta Q_k.
%	\label{eq:realteil_skalar}
%\end{equation}
%Bemerkenswert ist, dass der Widerstandsanteil $R^{(\ell)}_{kk}$ ausschließlich zum
%Imaginärteil beiträgt und sich somit vollständig herauskürzt. Einsetzen
%in~\eqref{eq:betragsquadrat_identitaet} liefert die zentrale
%Sensitivitätsbeziehung:
%\begin{equation}
%	\Delta(|V_k|^2) \;=\; 2\,X^{(\ell)}_{kk}\cdot\Delta Q_k.
%	\label{eq:sensitivity_skalar}
%\end{equation}

%----------------------------------------------------------------------
%\paragraph{Ergebnis: Q-Korrekturformel.}
%Um die Spannung vom berechneten Wert $|v^{(\ell)}_k|$ auf den Sollwert
%$|V^{\spec}_k|$ zu korrigieren, wird $\Delta(|V_k|^2) = |V^{\spec}_k|^2 -
%|v^{(\ell)}_k|^2$ gefordert. Auflösen von~\eqref{eq:sensitivity_skalar} nach
%$\Delta Q_k$ ergibt:
\begin{equation}
	\boxed{\;\Delta Q^{(\ell)}_k \;=\; \frac{|V^{\spec}_k|^2 - |V^{(\ell)}_k|^2}{2\cdot X^{(\ell)}_{kk}}, \;}
	\label{eq:methode_a_dQ_scalar}
\end{equation}
%wobei $X^{(\ell)}_{kk} = \im\bigl(\bigl[(\vect{B}^{(\ell)})^{-1}\bigr]_{kk}\bigr)$
%die Thévenin-Reaktanz am Knoten~$k$ ist. 

hergeleitet.
Dieser kann folglich als Korrekturformel verwendet werden, nach dem im inneren TPF ein Spannungsvektor berechnet wird.


%======================================================================
\subsection{Berücksichtigung der PV-PV-Kopplung}

Jedoch wird in \cref{sec:Konvergenz der äußeren Schleife: Entkoppelte Sensitivität} der Ergebnisse gezeigt, dass diese reine Thévenin-Formulierung der Sensitivität einen großen Nachteil hat.
Die zuvor vorgeschlagene Herleitung~\eqref{eq:thevenin_skalar}--\eqref{eq:methode_a_dQ_scalar} behandelt jeden PV-Knoten einzeln.
Dies hat zur Folge, dass bei der ausschließlichen Verwendung von den Diagonalelementen von~$Z^{(\ell)}_{B,kk}$ die Kopplung von zwei verschiedenen PV-Knoten unterschlagen wird.
Jedoch kann diese Annahme nur gelten, wenn mehrere PV-Knoten keinen Einfluss aufeinander haben.
Liegen mehrerer solcher Knoten im Netz vor, können sie sich Leitungen teile. 
Folglich kann eine Blindleistungsänderung an einem Knoten die Spannung eines anderen PV-Knotens beeinflussen.
Dieses Verhalten lässt sich bei Erhöhung von PV-Knoten eines radialen Netzes in \cref{sec:Konvergenz der äußeren Schleife: Entkoppelte Sensitivität} unter Konvergenzverlust beobachten.\\

Deswegen wird hier alternative modifizierte Sensitivität für die Blindleistungskorrektur vorgeschlagen, welche die elektrische Kopplung mehrerer PV-Knoten berücksichtigt. 
Da die Kopplungen zwischen Knoten werden im Netzmodell (\cref{subsec:admittanz}) durch die Nebendiagonalelemente der Impedanzmatrix modelliert. 
Dies motiviert eine konsequente Erweiterung der skalaren Formulierung der Sensitivität \eqref{eq:methode_a_dQ_scalar} in eine Matrixform, um die gesamte Impedanzmatrix verwenden zu können.
Die Skalare $\Delta Q_k$ und $\Delta V_{k}$ werden durch die korrespondierenden Vektorgrößern $\Delta\mathbf{Q}$,
$\Delta\mathbf{V}_{\mathrm{pv}}$ erweitert, wobei der $k$-te PV-Knoten in der $k$-ten Zeile der Vektoren steht.
Für die Impedanz wird aus der Geamtmatrix $\ZB$ gebaut, indem alle PV-PV-Einträge zu einem Block $\Delta\mathbf{V}_\mathrm{pv}$ zusammengetragen werden.
Die Nebendiagonalelemente dieser Matrix beschreiben genau den Einfluss der PV-Knoten auf einander.\\

Die einzelnen Rechenschritte der Herleitung bleiben die selben.
aus der skalaren entkoppelten Gleichung 

\begin{equation}
	\;\Delta Q^{(\ell)}_k \;=\; \frac{|V^{\spec}_k|^2 - |V^{(\ell)}_k|^2}{2\cdot X^{(\ell)}_{kk}}
	\label{eq:skale}
\end{equation} 

wird folglich 

\begin{equation}
	\boxed{\;\Delta\vect{Q}^{(\ell)} \;=\; \tfrac{1}{2}\,\bigl(\vect{X}^{(\ell)}_{pp}\bigr)^{-1} \bigl(|\vect{V}^{\spec}_{pv}|^2 - |\vect{v}^{(\ell)}_{pv}|^2\bigr) \;}
	\label{eq:methode_a_dQ_coupled}
\end{equation}

als gekoppelte Blindleistungskorrektur aufgefasst.
Werden in dieser bloß die Hauptdiagonalelemente betrachtet, entsprechen die Ausdrücke der entkoppelten Formulierung in \eqref{eq:skale}.
Es wird ebenfalls in \cref{sec:Konvergenz der äußeren Schleife: Entkoppelte Sensitivität} empirisch gezeigt, dass diese Formulierung deutlich robuster zu einer Konvergenz der Berechnung führt, als die entkoppelte Variante.


\subsection{Tensorformulierung der äußeren Blindleistungskorrektur}

Die hergelittenen Formulierungen für die Blindleistungskorrektur beziehen sich zunächst auf einen einzigen Lastfluss.
In diesem Abschnitt sollen die beiden Varianten der $Q$-Korrektur nun in eine Tensorformulierung überführt werden, um der Anforderung in  \cref{sec:pv_auswirkung} gerecht zu werden.
Dabei soll es möglich sein diese Korrektur auch für $\tau$-Lastflüsse gleichzeitig ausführen zu können.
Bei mehreren Lastszenarien liegen die berechneten Größen als mehrdimensionale Tensoren vor, wie in der Dense-Formulierung bereits eingeführt wurde.
Die Spannungen 
\begin{equation}
	|\dot{{V}}_{\mathrm{pv}}^{\spec}|^2,
	\qquad
	|\dot{{v}}_{\mathrm{pv}}^{(\ell)}|^2
	\qquad \in
	\mathbb{R}^{n_{\mathrm{pv}}\times\tau}
	\label{eq:dense_pv_voltage}
\end{equation}
liegen dabei bereits in tensorisierter Form durch die Berechnung vor und dienen als Ausgangspunkt für die folgenden Formulierungen.
Die Blindleistung

\begin{equation}
	\Delta\dot{{Q}}^{(\ell)}
	\in
	\mathbb{R}^{n_{\mathrm{pv}}\times\tau}
	\label{eq:dense_dQ}
\end{equation}

soll ebenfalls mehrdimensional vorliegen um diese einfach in den TPF übernehmen zu können.
Die Spalten beherbergen dabei die einzelnen Lastflüsse und die Zeilen wie zuvor die PV-Busse.
Für das Spannungsresiduum folgt eine einfache spaltenweise Subtraktion der Elemente

\begin{equation}
	\dot{{R}}^{(\ell)}
	=
	|\dot{{V}}_{\mathrm{pv}}^{\spec}|^2
	-
	|\dot{{v}}_{\mathrm{pv}}^{(\ell)}|^2
	=
	\begin{bmatrix}
		r_{1,1}^{(\ell)} & r_{1,2}^{(\ell)} & \cdots & r_{1,\tau}^{(\ell)}\\
		r_{2,1}^{(\ell)} & r_{2,2}^{(\ell)} & \cdots & r_{2,\tau}^{(\ell)}\\
		\vdots & \vdots & \ddots & \vdots\\
		r_{n_{\mathrm{pv}},1}^{(\ell)}
		&
		r_{n_{\mathrm{pv}},2}^{(\ell)}
		&
		\cdots
		&
		r_{n_{\mathrm{pv}},\tau}^{(\ell)}
	\end{bmatrix}
	\in
	\mathbb{R}^{n_{\mathrm{pv}}\times\tau},
	\label{eq:dense_pv_residual}
\end{equation}
wobei ebenfalls eine zweidimensionale Matrix folgt.
Da für die betrachtete Formulierung eine konstante Netztopologie
angenommen wird, ist die PV-PV-Sensitivitätsmatrix $\vect{X}_{pp}$
für alle Szenarien identisch. Die Sensitivität muss daher nicht für jedes
einzelne Szenario wiederholt werden. Stattdessen kann dieselbe Matrix durch
eine Matrixmultiplikation auf sämtliche Spalten des Spannungsresiduums
angewendet werden.\\


Für die entkoppelte Formulierung werden die
Hauptdiagonalelemente der PV-PV-Reaktanzmatrix berücksichtigt. Die
entsprechende Diagonalmatrix ist durch
\begin{equation}
	\vect{X}_{pp}^{\mathrm{diag}}
	=
	\operatorname{diag}
	\left(
	\vect{X}_{pp}
	\right)
	\in
	\mathbb{R}^{n_{\mathrm{pv}}\times n_{\mathrm{pv}}}
	\label{eq:dense_Xpp_diag}
\end{equation}
gegeben.
Die skalare Blindleistungskorrektur
\begin{equation}
	\Delta Q_k^{(\ell)}
	=
	\frac{
		|V_k^{\spec}|^2
		-
		|v_k^{(\ell)}|^2
	}{
		2X_{kk}
	}
	\label{eq:dense_scalar_decoupled}
\end{equation}
kann damit für sämtliche PV-Knoten und Szenarien gleichzeitig als
Matrixoperation 
\begin{equation}
	\boxed{
		\Delta\dot{\vect{Q}}_{\mathrm{dec}}^{(\ell)}
		=
		\frac{1}{2}
		\left(
		\vect{X}_{pp}^{\mathrm{diag}}
		\right)^{-1}
		\cdot
		\left(
		|\dot{\vect{V}}_{\mathrm{pv}}^{\spec}|^2
		-
		|\dot{\vect{v}}_{\mathrm{pv}}^{(\ell)}|^2
		\right)
	}
	\label{eq:dense_decoupled_Q}
\end{equation}

formuliert werden.
Die Elemente können dabei Element-weise mit einander verrechnet werden.\\

Unter Berücksichtigung der PV-Kopplungen 
werden zusätzlich die Nebendiagonalelemente der
PV-PV-Reaktanzmatrix~$\vect{X}_{pp}$ verwendet.
Da die restlichen Bedingungen die selben sind wie in der entkoppelten Variante, folgt für das Residuum
\begin{align}
	\dot{{R}}^{(\ell)}
	&=
	2\vect{X}_{pp}
	\cdot
	\Delta\dot{{Q}}^{(\ell)}\\
	&=
	2
	\begin{bmatrix}
		X_{11} & X_{12} & \cdots & X_{1n_{\mathrm{pv}}}\\
		X_{21} & X_{22} & \cdots & X_{2n_{\mathrm{pv}}}\\
		\vdots & \vdots & \ddots & \vdots\\
		X_{n_{\mathrm{pv}}1}
		&
		X_{n_{\mathrm{pv}}2}
		&
		\cdots
		&
		X_{n_{\mathrm{pv}}n_{\mathrm{pv}}}
	\end{bmatrix}
	\begin{bmatrix}
		\Delta Q_{1,1}^{(\ell)}
		&
		\Delta Q_{1,2}^{(\ell)}
		&
		\cdots
		&
		\Delta Q_{1,\tau}^{(\ell)}
		\\
		\Delta Q_{2,1}^{(\ell)}
		&
		\Delta Q_{2,2}^{(\ell)}
		&
		\cdots
		&
		\Delta Q_{2,\tau}^{(\ell)}
		\\
		\vdots
		&
		\vdots
		&
		\ddots
		&
		\vdots
		\\
		\Delta Q_{n_{\mathrm{pv}},1}^{(\ell)}
		&
		\Delta Q_{n_{\mathrm{pv}},2}^{(\ell)}
		&
		\cdots
		&
		\Delta Q_{n_{\mathrm{pv}},\tau}^{(\ell)}
	\end{bmatrix}
	\label{eq:dense_coupled_sensitivity}
\end{align}

Die gekoppelte Blindleistungskorrektur in Tensorform ist folglich durch
\begin{equation}
	\boxed{
		\Delta\dot{{Q}}^{(\ell)}
		=
		\frac{1}{2}
		\vect{X}_{pp}^{-1}
		\cdot
		\left(
		|\dot{\vect{V}}_{\mathrm{pv}}^{\spec}|^2
		-
		|\dot{\vect{v}}_{\mathrm{pv}}^{(\ell)}|^2
		\right)
	}.
	\label{eq:dense_coupled_Q}
\end{equation}
gegeben



%
%Die Herleitung~\eqref{eq:thevenin_skalar}--\eqref{eq:methode_a_dQ_scalar}
%behandelt jeden PV-Knoten für sich und verwendet dabei nur das
%Diagonalelement~$Z^{(\ell)}_{B,kk}$. Damit wird angenommen, dass
%eine Q-Änderung am Knoten~$k$ ausschließlich die Spannung an genau diesem Knoten
%beeinflusst. Tatsächlich teilen sich die PV-Knoten in einem realen Netz jedoch
%gemeinsame Leitungen, sodass eine Blindleistungsänderung am Knoten~$j$ auch die
%Spannung am Knoten~$k$ verschiebt. Diese Kopplung lässt sich erfassen, indem man
%die skalare Rechnung konsequent in Matrixnotation überführt. An die Stelle der
%Skalare treten Vektoren über alle PV-Knoten ($\Delta\vect{Q}$,
%$\Delta\vect{V}_{pv}$), und das Diagonalelement~$Z^{(\ell)}_{B,kk}$ wird durch den
%gesamten PV--PV-Block~$\vect{Z}^{(\ell)}_{pp}$ der Bus-Impedanzmatrix ersetzt.
%Dessen Off-Diagonalelemente~$(k\neq j)$ tragen gerade die zuvor vernachlässigte
%Kopplung.
%
%Die einzelnen Schritte bleiben unverändert, werden aber vektoriell gelesen. Aus
%der knotenweisen Stromänderung~\eqref{eq:strom_linear} wird
%$\Delta\vect{I}_{pv} = -\jj\,\vect{D}_V^{-1}\,\Delta\vect{Q}$ mit der
%Diagonalmatrix $\vect{D}_V = \diag((v^{(\ell)}_k)^*)$, und aus dem
%Thévenin-Schritt~\eqref{eq:spannung_skalar} wird
%$\Delta\vect{V}_{pv} = \vect{Z}^{(\ell)}_{pp}\,\Delta\vect{I}_{pv}$. Wertet man
%$\Delta(|V_k|^2) = 2\,\re(V_k^*\Delta V_k)$ komponentenweise aus und nutzt für
%Verteilnetze die Näherungen ähnlicher Spannungsbeträge $(|v^{(\ell)}_k| \approx
%|v^{(\ell)}_j|)$ und kleiner Winkeldifferenzen $(\theta_j - \theta_k \approx 0)$,
%so kürzen sich die Widerstandsanteile wie im skalaren Fall heraus. Es verbleibt die
%matrixwertige Sensitivitätsbeziehung
%\begin{equation}
%	\Delta\bigl(|\vect{V}_{pv}|^2\bigr) \;=\; 2\,\vect{X}^{(\ell)}_{pp}\,\Delta\vect{Q},
%	\label{eq:sensitivity_matrix}
%\end{equation}
%die direkte Verallgemeinerung von~\eqref{eq:sensitivity_skalar}: Aus dem Skalar
%$2X^{(\ell)}_{kk}$ ist die Matrix $2\vect{X}^{(\ell)}_{pp}$ geworden, deren
%Diagonale genau die skalaren Sensitivitäten enthält. Auflösen
%nach~$\Delta\vect{Q}$ ergibt die
%gekoppelte Q-Korrekturformel
%\begin{equation}
%	\boxed{\;\Delta\vect{Q}^{(\ell)} \;=\; \tfrac{1}{2}\,\bigl(\vect{X}^{(\ell)}_{pp}\bigr)^{-1} \bigl(|\vect{V}^{\spec}_{pv}|^2 - |\vect{v}^{(\ell)}_{pv}|^2\bigr) \;}
%	\label{eq:methode_a_dQ_coupled}
%\end{equation}
%Vernachlässigt man die Off-Diagonalen von $\vect{X}^{(\ell)}_{pp}$, zerfällt
%\eqref{eq:methode_a_dQ_coupled} knotenweise wieder in die skalare Formel
%\eqref{eq:methode_a_dQ_scalar}. Beide Varianten sind gleichermaßen
%tensorisierbar, da $\vect{X}^{(\ell)}_{pp}$ über alle Iterationen und Szenarien





%konstant ist, lässt sich $(\vect{X}^{(\ell)}_{pp})^{-1}$ einmalig vorab berechnen,
%und das Q-Update über alle $\tau$ Szenarien reduziert sich auf ein
%Matrix-Matrix-Produkt -- denselben Operationstyp, der auch den TPF-Kern
%\eqref{eq:tpf_iteration} beschleunigt. Der Unterschied ist daher nicht die
%Parallelisierbarkeit, sondern der Aufwand pro äußerer Iteration (elementweises
%gegenüber dichtem Produkt) sowie die Frage, ob die Kopplung zwischen den
%PV-Knoten in der Suchrichtung berücksichtigt wird. Zudem läuft das Q-Update in
%der äußeren Schleife unabhängig vom tensorisierten inneren TPF ab, sodass es den
%Parallelisierungsvorteil (F2) ohnehin nicht dominiert.


%
%
%
%
%
%Für die Q-Korrektur in Schritt~5 des Ablaufs benötigen wir eine Vorschrift, die
%uns sagt, \emph{um wie viel} die Blindleistung an jedem PV-Knoten geändert werden
%muss, damit die Spannungsbeträge ihre Sollwerte erreichen. In diesem Abschnitt
%leiten wir diese Vorschrift Schritt für Schritt her. Der rote Faden ist dabei das
%\emph{Newton-Verfahren}: Wir formulieren die PV-Bedingung als Nullstellenproblem,
%wenden darauf den (mehrdimensionalen) Newton-Schritt an und berechnen die dafür
%benötigte Ableitung (die \emph{Sensitivität}) analytisch über ein
%Thévenin-Ersatzmodell des Netzes.
%
%%----------------------------------------------------------------------
%\paragraph{Notation.}
%Bevor wir loslegen, führen wir kurz die wichtigsten Größen ein. Ein Netz mit
%PV-Knoten hat $n_{pv}$ solcher Knoten. Für jeden davon suchen wir die
%\emph{Blindleistung}, denn genau die ist an einem PV-Knoten unbekannt (Wirkleistung
%und Spannungsbetrag sind ja vorgegeben). Wir schreiben diese unbekannten
%Blindleistungen aller PV-Knoten in einen Vektor $\vect{q}$. Analog fassen wir die
%zugehörigen \emph{Spannungszeiger} zum Vektor $\vect{V}_{pv}$ und die
%\emph{Strominjektionen} zum Vektor $\vect{I}_{pv}$ zusammen -- jeder Eintrag
%gehört zu genau einem PV-Knoten.
%
%Die entscheidende Netzgröße ist die \emph{Bus-Impedanzmatrix}
%$\vect{Z}_B^{(\ell)} = \bigl(\vect{B}^{(\ell)}\bigr)^{-1}$. Anschaulich beschreibt
%ihr Eintrag in Zeile~$k$ und Spalte~$j$, um wie viel die Spannung am Knoten~$k$
%ansteigt, wenn man am Knoten~$j$ einen Strom einspeist. Der obere Index~$(\ell)$
%erinnert daran, dass sie sich mit der äußeren Iteration ändern kann. Da uns nur
%die PV-Knoten interessieren, betrachten wir den Ausschnitt dieser Matrix, der nur
%die Zeilen und Spalten der PV-Knoten enthält, und nennen ihn $\vect{Z}_{pp}^{(\ell)}$
%(die ``PV-zu-PV''-Untermatrix). Von dieser Untermatrix brauchen wir am Ende nur
%den Imaginärteil, also den \emph{Reaktanz-Anteil}; ihn bezeichnen wir mit
%$\vect{X}_{pp}^{(\ell)} = \im\bigl(\vect{Z}_{pp}^{(\ell)}\bigr)$. Warum
%ausgerechnet dieser Reaktanz-Anteil die zentrale Rolle spielt, wird sich im Lauf
%der Herleitung zeigen.
%%----------------------------------------------------------------------
%\paragraph{Schritt 0: Das PV-Problem als Nullstellensuche.}
%An jedem PV-Knoten $k \in \mathcal{P}$ ist der Wirkleistungssollwert
%$P_k^{\spec}$ und der Spannungsbetrags-Sollwert $|V_k^{\spec}|$ bekannt, die
%Blindleistung $q_k$ jedoch \emph{unbekannt}. Physikalisch existiert genau der
%Wert $q_k$, der dafür sorgt, dass die tatsächliche Spannung $|V_k|$ ihren Sollwert
%trifft. Wir suchen also den Vektor $\vect{q}$, für den das \emph{Residuum}
%\begin{equation}
%	\vect{g}(\vect{q}) \;=\; |\vect{V}^{\spec}|^2 - |\vect{V}(\vect{q})|^2
%	\;\in\; \mathbb{R}^{n_{pv}}
%	\label{eq:residuum_intro}
%\end{equation}
%verschwindet, d.\,h.\ $\vect{g}(\vect{q}^\star) = \vect{0}$. Die $k$-te
%Komponente $g_k$ beschreibt, wie weit der Knoten~$k$ noch von seinem
%Spannungsziel entfernt ist. Wir arbeiten hier bewusst mit dem
%\emph{Betragsquadrat} $|V_k|^2$ statt mit $|V_k|$, weil sich $|V_k|^2$ -- wie wir
%gleich sehen -- besonders einfach ableiten lässt.
%
%Wichtig für das Verständnis: $|\vect{V}(\vect{q})|^2$ ist \emph{keine} einfache
%Formel, sondern das Ergebnis eines vollständigen Lastflusses. Man gibt einen
%Blindleistungsvektor $\vect{q}$ vor, lässt die innere TPF-Schleife bis zur
%Konvergenz laufen und liest die Spannungen ab. Die Funktion $\vect{g}$ ist also
%eine ``Black-Box'' in $\vect{q}$ -- genau die Situation, für die das
%Newton-Verfahren entwickelt wurde.
%
%%----------------------------------------------------------------------
%\paragraph{Schritt 1: Der Newton-Schritt (vom Skalar zum Vektor).}
%Zur Erinnerung der eindimensionale Fall (\emph{ein} PV-Knoten): Um eine
%Nullstelle von $g(q)$ zu finden, iteriert das Newton-Verfahren
%\begin{equation}
%	q^{(\ell+1)} \;=\; q^{(\ell)} - \frac{g\bigl(q^{(\ell)}\bigr)}{g'\bigl(q^{(\ell)}\bigr)}.
%	\label{eq:newton_skalar}
%\end{equation}
%Die Ableitung $g'(q) = -\,\partial|V|^2/\partial q$ ist die \emph{Sensitivität}:
%Sie gibt an, um wie viel sich das Spannungsbetragsquadrat ändert, wenn man die
%Blindleistung ändert.
%
%Bei mehreren PV-Knoten wird $\vect{q}$ ein Vektor, $\vect{g}$ ein Vektor und die
%Ableitung eine \emph{Matrix} -- die Jacobi-Matrix. Wir definieren die
%\emph{Sensitivitätsmatrix}
%\begin{equation}
%	\vect{H}_{kj} \;=\; \frac{\partial |V_k|^2}{\partial q_j},
%	\qquad k, j \in \mathcal{P},
%	\label{eq:H_def_intro}
%\end{equation}
%sodass $\partial\vect{g}/\partial\vect{q} = -\vect{H}$ gilt. Der mehrdimensionale
%Newton-Schritt lautet damit
%\begin{equation}
%	\vect{q}^{(\ell+1)}
%	\;=\; \vect{q}^{(\ell)} + \vect{H}^{-1}\,
%	\bigl(|\vect{V}^{\spec}|^2 - |\vect{v}^{(\ell)}_{pv}|^2\bigr).
%	\label{eq:newton_vektor}
%\end{equation}
%\textbf{Dies ist bereits die vollständige Motivation der äußeren Schleife}: Sie
%ist nichts anderes als ein Newton-Verfahren auf den PV-Spannungsbedingungen. Es
%bleibt nur die Frage, wie wir die Sensitivitätsmatrix $\vect{H}$ berechnen.
%
%\medskip
%\emph{Warum die Kopplung zwingend auftritt.} Das Off-Diagonalelement
%$\vect{H}_{kj}$ mit $k \neq j$ beantwortet die Frage: ``Wie stark ändert sich die
%Spannung an Knoten~$k$, wenn ich die Blindleistung an Knoten~$j$ ändere?'' Dieser
%Wert ist ungleich null, weil sich die PV-Knoten gemeinsame Netzimpedanzen teilen.
%Vernachlässigt man ihn (die \emph{entkoppelte} bzw.\ diagonale Näherung), so tut
%man so, als sei jeder PV-Knoten elektrisch isoliert. Genau diese Kopplung
%berücksichtigen wir im Folgenden.
%
%%----------------------------------------------------------------------
%\paragraph{Schritt 2: Das Netz als lineare Abbildung (multivariates Thévenin).}
%Um $\vect{H}$ analytisch zu bestimmen, brauchen wir einen Ausdruck dafür, wie die
%Spannungen von den Strominjektionen abhängen. Ausgangspunkt ist die Netzgleichung
%$\vect{B}^{(\ell)}\vect{V} + \vect{Y}_{ds}\vect{v}_s = \vect{I}$. Auflösen nach
%$\vect{V}$ liefert
%\begin{equation}
%	\vect{V} \;=\; \underbrace{-\,\vect{Z}_B^{(\ell)}\,\vect{Y}_{ds}\,\vect{v}_s}_{\vect{V}^{\text{open}}}
%	\;+\; \vect{Z}_B^{(\ell)}\,\vect{I},
%	\qquad
%	I_k = \frac{s_k^*}{V_k^*}.
%	\label{eq:thevenin_intro}
%\end{equation}
%Dies ist die Mehrknoten-Verallgemeinerung des Thévenin-Theorems:
%$\vect{V}^{\text{open}}$ ist die Leerlaufspannung (kein Strom injiziert), und
%$\vect{Z}_B^{(\ell)}[k,j]$ beschreibt den Spannungsanstieg an Knoten~$k$ pro
%Einheit Stromeinspeisung an Knoten~$j$. Einschränkung auf die PV-Zeilen ergibt
%\begin{equation}
%	\vect{V}_{pv}
%	\;=\; \vect{V}^{\text{open}}_{pv}
%	\;+\; \vect{Z}_{pp}^{(\ell)}\,\vect{I}_{pv}
%	\;+\; \underbrace{\sum_{j \notin \mathcal{P}} \vect{Z}_B^{(\ell)}[\mathcal{P}, j]\,I_j}_{\text{PQ-Anteil, konstant}}.
%	\label{eq:thevenin_pv_intro}
%\end{equation}
%Der letzte Term hängt nur von den PQ-Strömen ab und bleibt während einer reinen
%Q-Variation an den PV-Knoten in erster Näherung konstant. Entscheidend ist:
%\emph{$\vect{Z}_{pp}^{(\ell)}$ ist im Allgemeinen keine Diagonalmatrix}. Ein
%Off-Diagonalelement $\vect{Z}_{pp}^{(\ell)}[k,j]$ entspricht in radialen Netzen
%der Summe der Serienimpedanzen entlang des gemeinsamen Vorfahrenpfads von der
%Wurzel (Slack) zu $k$ bzw.\ $j$. Zwei PV-Knoten im selben Ast sind daher stark
%gekoppelt.
%
%%----------------------------------------------------------------------
%\paragraph{Schritt 3: Die Sensitivität über die Kettenregel.}
%Die gesuchte Ableitung $\partial|V_k|^2/\partial q_j$ setzt sich aus drei
%einfachen Teilableitungen zusammen (Kettenregel):
%\begin{equation}
%	\frac{\partial |V_k|^2}{\partial q_j}
%	\;=\;
%	\underbrace{\frac{\partial |V_k|^2}{\partial V_k}}_{\text{(a) Betrag}}
%	\cdot
%	\underbrace{\frac{\partial V_k}{\partial I_j}}_{\text{(b) Netz}}
%	\cdot
%	\underbrace{\frac{\partial I_j}{\partial q_j}}_{\text{(c) Einspeisung}}.
%	\label{eq:kettenregel_intro}
%\end{equation}
%
%\subparagraph{Teil (c): Strom als Funktion der Blindleistung.}
%Am PV-Knoten $j$ gilt in Erzeugerkonvention $I_j = (P_j + \jj q_j)^*/V_j^*$. Bei
%festgehaltener Wirkleistung $P_j$ und festgehaltener Spannung $V_j$ (Linearisierung
%um den aktuellen Arbeitspunkt $v_j^{(\ell)}$) folgt
%\begin{equation}
%	\Delta I_j = \frac{-\jj\,\Delta q_j}{V_j^*}\bigg|_{V_j = v_j^{(\ell)}}
%	\quad\Longrightarrow\quad
%	\Delta\vect{I}_{pv} = -\jj\,\vect{D}_{V^*}^{-1}\,\Delta\vect{q},
%	\label{eq:teil_c}
%\end{equation}
%mit der Diagonalmatrix
%$\vect{D}_{V^*} = \diag\bigl(v_{k_1}^{(\ell)*}, \ldots, v_{k_{n_{pv}}}^{(\ell)*}\bigr)$.
%
%\subparagraph{Teil (b): Spannung als Funktion des Stroms.}
%Aus \eqref{eq:thevenin_pv_intro} sind alle Terme außer
%$\vect{Z}_{pp}^{(\ell)}\vect{I}_{pv}$ in der Linearisierung konstant, also
%\begin{equation}
%	\Delta\vect{V}_{pv}
%	\;=\; \vect{Z}_{pp}^{(\ell)}\,\Delta\vect{I}_{pv}
%	\;=\; -\jj\,\vect{Z}_{pp}^{(\ell)}\,\vect{D}_{V^*}^{-1}\,\Delta\vect{q}.
%	\label{eq:teil_b}
%\end{equation}
%\emph{Hier steckt die gesamte Kopplung}: Über die Off-Diagonalen von
%$\vect{Z}_{pp}^{(\ell)}$ wirkt eine Stromänderung an $j$ auf die Spannung an $k$.
%Komponentenweise:
%\begin{equation}
%	\Delta V_k
%	\;=\; \sum_{j \in \mathcal{P}}
%	\vect{Z}_{pp}^{(\ell)}[k, j]\cdot\frac{-\jj\,\Delta q_j}{v_j^{(\ell)*}}.
%	\label{eq:delta_V_komp}
%\end{equation}
%
%\subparagraph{Teil (a): Vom komplexen Zeiger zum Betragsquadrat.}
%Für das Betragsquadrat gilt exakt die Identität
%\begin{equation}
%	\Delta\bigl(|V_k|^2\bigr) \;=\; 2\,\re\!\bigl(V_k^*\,\Delta V_k\bigr).
%	\label{eq:teil_a}
%\end{equation}
%Der Faktor~$2$ stammt genau aus dieser Ableitung des Quadrats. Setzt man
%\eqref{eq:delta_V_komp} ein und schreibt $V_k = |v_k^{(\ell)}|\,e^{\jj\theta_k}$,
%so entsteht der Faktor
%\begin{equation}
%	V_k^* \cdot \frac{1}{v_j^{(\ell)*}}
%	\;=\; \frac{|v_k^{(\ell)}|}{|v_j^{(\ell)}|}\,e^{\jj(\theta_j - \theta_k)}.
%	\label{eq:faktor}
%\end{equation}
%
%%----------------------------------------------------------------------
%\paragraph{Schritt 4: Vereinfachung für Verteilnetze.}
%In einem normal betriebenen radialen Verteilnetz gelten zwei Näherungen:
%
%\begin{center}
%	\begin{tabular}{@{}lll@{}}
%		\toprule
%		\textbf{Annahme} & \textbf{Bedeutung} & \textbf{Begründung} \\
%		\midrule
%		$|v_k^{(\ell)}| \approx |v_j^{(\ell)}|$ & alle PV-Spannungen nahe $1$~p.u. & spannungsgeregelte Knoten \\
%		$\theta_j - \theta_k \approx 0$ & kleine Winkeldifferenzen & kurze Leitungen, radiale Struktur \\
%		\bottomrule
%	\end{tabular}
%\end{center}
%
%Unter diesen Annahmen reduziert sich der Faktor \eqref{eq:faktor} zu~$1$. Mit der
%Zerlegung $\vect{Z}_{pp}^{(\ell)}[k,j] = R_{kj}^{(\ell)} + \jj X_{kj}^{(\ell)}$
%folgt aus \eqref{eq:delta_V_komp} und \eqref{eq:faktor}
%\begin{equation}
%	V_k^*\,\Delta V_k
%	\;\approx\; \sum_{j \in \mathcal{P}}
%	\bigl(X_{kj}^{(\ell)} - \jj R_{kj}^{(\ell)}\bigr)\,\Delta q_j.
%	\label{eq:VstarDeltaV}
%\end{equation}
%Der Realteil liefert dann
%\begin{equation}
%	\re\!\bigl(V_k^*\,\Delta V_k\bigr)
%	\;=\; \sum_{j \in \mathcal{P}} X_{kj}^{(\ell)}\,\Delta q_j
%	\;=\; \bigl[\vect{X}_{pp}^{(\ell)}\,\Delta\vect{q}\bigr]_k.
%	\label{eq:realteil}
%\end{equation}
%Bemerkenswert: Sämtliche Widerstandsanteile $R_{kj}^{(\ell)}$ entfallen -- auch
%die Off-Diagonalen. Die Spannungsbetragsänderung wird allein durch die
%Reaktanzen bestimmt.
%
%%----------------------------------------------------------------------
%\paragraph{Schritt 5: Zentrale Sensitivitätsbeziehung.}
%Zusammen mit dem Faktor~$2$ aus \eqref{eq:teil_a} erhalten wir die kompakte
%Matrixform für den gesamten PV-Vektor:
%\begin{equation}
%	\boxed{\;\Delta|\vect{V}_{pv}|^2 \;=\; 2\,\vect{X}_{pp}^{(\ell)}\,\Delta\vect{q}\;}
%	\label{eq:sensitivity_matrix}
%\end{equation}
%Damit ist die gesuchte Jacobi-Matrix aus \eqref{eq:H_def_intro} in erster Ordnung
%gefunden:
%\begin{equation}
%	\vect{H} \;\approx\; 2\,\vect{X}_{pp}^{(\ell)}.
%	\label{eq:H_gleich_2X}
%\end{equation}
%Das multivariate Thévenin-Modell hat also genau die Ableitung geliefert, die das
%Newton-Verfahren benötigt. Für $n_{pv} = 1$ ist $\vect{X}_{pp}^{(\ell)}$ eine
%$1\times1$-Matrix mit dem einzigen Eintrag $X_{kk}^{(\ell)}$, und
%\eqref{eq:sensitivity_matrix} reduziert sich auf die klassische skalare Beziehung
%$\Delta(|V_k|^2) = 2 X_{kk}^{(\ell)}\,\Delta q_k$.
%
%%----------------------------------------------------------------------
%\paragraph{Ergebnis: Die gekoppelte Q-Korrekturformel.}
%Einsetzen von \eqref{eq:H_gleich_2X} in den Newton-Schritt
%\eqref{eq:newton_vektor} und Auflösen nach $\Delta\vect{q}$ liefert die gesuchte
%Korrektur:
%\begin{equation}
%	\boxed{\;
%		\Delta\vect{q}^{(\ell)}
%		\;=\; \tfrac{1}{2}\,\bigl(\vect{X}_{pp}^{(\ell)}\bigr)^{-1}
%		\bigl(|\vect{V}^{\spec}|^2 - |\vect{v}^{(\ell)}_{pv}|^2\bigr).
%		\;}
%	\label{eq:methode_a_dQ_coupled}
%\end{equation}
%Diese \emph{gekoppelte} Formel löst das gesamte PV-System simultan: Jede
%Blindleistungsänderung $\Delta q_j$ berücksichtigt über $\vect{X}_{pp}^{(\ell)}$
%die Rückwirkung auf alle anderen PV-Knoten. Sie ersetzt die \emph{entkoppelte}
%Näherung
%\begin{equation}
%	\Delta q_k^{(\ell)}
%	\;=\; \frac{|V_k^{\spec}|^2 - |v_k^{(\ell)}|^2}{2\,X_{kk}^{(\ell)}}
%	\qquad\text{(nur Diagonale, PV-Knoten isoliert betrachtet)},
%	\label{eq:methode_a_dQ_scalar}
%\end{equation}
%die man aus \eqref{eq:methode_a_dQ_coupled} erhält, wenn man alle Off-Diagonalen
%von $\vect{X}_{pp}^{(\ell)}$ streicht. Die entkoppelte Formel ist somit ein
%Spezialfall der gekoppelten.
%
%%----------------------------------------------------------------------
%\paragraph{Warum der gekoppelte Ansatz so schnell konvergiert.}
%Der Nutzen des korrekten Jacobians zeigt sich direkt im Konvergenzverhalten. Die
%Jacobi-Matrix der Fixpunktabbildung der äußeren Schleife ist (mit
%Dämpfungsfaktor~$\omega$ und $\vect{A} = 2\vect{X}_{pp}$)
%\begin{equation}
%	\vect{J}_G = \vect{I} - \omega\,\vect{A}^{-1}\vect{H}.
%\end{equation}
%Da unser Modell $\vect{H} \approx \vect{A}$ liefert, folgt
%$\vect{J}_G \approx (1-\omega)\,\vect{I}$ und damit $\rho(\vect{J}_G) \approx
%|1-\omega|$. Für $\omega = 1$ verschwindet die linearisierte Fehlerdämpfung: Die
%äußere Schleife konvergiert im linearen Regime in einer einzigen Iteration --
%\emph{unabhängig} von der Anzahl der PV-Knoten. Die entkoppelte Näherung
%\eqref{eq:methode_a_dQ_scalar} verwendet dagegen einen fehlerhaften Jacobian;
%ihre Off-Diagonalterme wachsen in radialen Bäumen mit der Zahl der PV-Knoten im
%selben Ast, sodass die Konvergenzbedingung oberhalb einer kritischen PV-Anzahl
%verletzt wird. Genau dieses Verhalten wird in \cref{ch:ergebnisse} numerisch
%untersucht.
%
%%----------------------------------------------------------------------
%\paragraph{Numerische Umsetzung.}
%Sind alle $\alpha_{Z, k} = 0$ an den PV-Knoten, so ist $\vect{B}^{(\ell)} =
%\vect{B}$ konstant über die äußere Iteration. Damit lässt sich
%$\vect{X}_{pp}$ (bzw.\ dessen LU-Faktorisierung oder Inverse) \emph{einmalig} vor
%Beginn der äußeren Schleife berechnen. Die Auswertung von
%\eqref{eq:methode_a_dQ_coupled} reduziert sich dann pro Iteration auf ein
%Matrix-Vektor-Produkt mit Kosten $\mathcal{O}(n_{pv}^2)$, was gegenüber den
%inneren FPI-Kosten $\mathcal{O}\bigl((b\varphi)^2\,\tau\bigr)$ vernachlässigbar
%ist ($n_{pv} \ll b\varphi$). Die Tensorstruktur bleibt erhalten:
%$(\vect{X}_{pp})^{-1}$ wird über alle $\tau$ Szenarien gebroadcastet.

%
%
%\subsection{Eigenschaften der Korrekturformel}
%\label{subsec:methode_a_eigenschaften}
%
%Die gekoppelte Formel~\eqref{eq:methode_a_dQ_coupled} besitzt mehrere
%bemerkenswerte Eigenschaften:
%
%\begin{enumerate}
%	\item \textbf{Widerstandsinvarianz.}
%	Wie im skalaren Fall entfallen sämtliche $R_{kj}^{(\ell)}$ vollständig;
%	\eqref{eq:methode_a_dQ_coupled} hängt ausschließlich von der
%	Reaktanz-Submatrix $\vect{X}_{pp}^{(\ell)}$ und den Spannungsbeträgen
%	ab, \emph{nicht} von den Winkeln $\theta_k$.
%	
%	\item \textbf{Physikalische Interpretation der Kopplung.}
%	Das Off-Diagonalelement $\vect{X}_{pp}^{(\ell)}[k, j]$ misst die
%	\emph{PV--PV-Sensitivität} $\partial|V_k|^2/\partial q_j$ und
%	entspricht in radialen Bäumen der Summe der Reaktanzen entlang des
%	\emph{gemeinsamen Vorfahrenpfads} von der Wurzel zu $k$ bzw.\ $j$:
%	\begin{equation}
%		\vect{X}_{pp}[k, j]
%		\;=\; \sum_{\text{line } \ell \,\in\, \mathrm{path}(0 \to k)\,\cap\,\mathrm{path}(0 \to j)}
%		x_\ell .
%		\label{eq:tree_impedance_interpretation}
%	\end{equation}
%	Zwei PV-Knoten im selben Ast eines Baums teilen fast ihren gesamten
%	Pfad und haben daher $\vect{X}_{pp}[k,j] \sim \vect{X}_{pp}[k,k]$.
%	Die entkoppelte Approximation vernachlässigt genau diesen Term.
%	
%	\item \textbf{Erhaltung der Struktur für $n_{pv} = 1$.}
%	Für einen einzelnen PV-Knoten wird~\eqref{eq:methode_a_dQ_coupled} zur
%	klassischen skalaren Formel~\eqref{eq:methode_a_dQ_scalar}. Die
%	entkoppelte Version ist ein Spezialfall der gekoppelten.
%	
%	\item \textbf{Abhängigkeit von $\vect{B}^{(\ell)}$.}
%	\begin{itemize}
%		\item Wenn $\alpha_{Z, k} = 0$ an allen PV-Knoten:
%		$\vect{X}_{pp}^{(\ell)} = \vect{X}_{pp}$ ist konstant; die
%		Faktorisierung wird einmalig vorab berechnet.
%		\item Wenn $\alpha_{Z, k} \neq 0$: $\vect{X}_{pp}^{(\ell)}$ ändert
%		sich mit $\vect{q}^{(\ell)}$; die Änderung ist meist gering und
%		$\vect{X}_{pp}^{(0)}$ ist eine gute Näherung.
%	\end{itemize}
%	
%	\item \textbf{Konditionierung.}
%	In radialen Netzen mit unterschiedlichen PV-Pfaden ist
%	$\vect{X}_{pp}$ symmetrisch positiv definit und wohlkonditioniert; die
%	Invertierung ist stabil. In pathologischen Fällen (mehrere PV-Knoten
%	an derselben elektrischen Position) empfiehlt sich eine Regularisierung
%	\begin{equation}
%		\vect{X}_{pp} \;\gets\; \vect{X}_{pp} + \eta_{\mathrm{reg}}\,\vect{I},
%		\qquad \eta_{\mathrm{reg}} \sim 10^{-10} .
%	\end{equation}
%	
%	\item \textbf{Lineare Konvergenz.}
%	\eqref{eq:methode_a_dQ_coupled} ist eine First-Order-Sensitivität
%	(Linearisierung um den Arbeitspunkt); die äußere Schleife konvergiert
%	linear. Der Konvergenzfaktor ist -- im Gegensatz zur entkoppelten
%	Variante -- \emph{unabhängig} von der Anzahl der PV-Knoten $n_{pv}$
%	(vgl.\ \S~\ref{subsec:methode_a_konvergenz}).
%\end{enumerate}

%\paragraph{Konvergenz der äußeren Schleife:}
%Die äußere Q-Schleife ist selbst eine Fixpunktiteration: Ausgehend vom Residuum
%$|V^{\spec}_k|^2 - |v^{(\ell)}_k|^2$ wird $\vect{Q}$ so lange korrigiert, bis die
%Spannungsbeträge ihre Sollwerte erreichen. Da der innere TPF pro äußerer
%Iteration exakt gelöst wird, bestimmt allein dieses Residuum die Konvergenz, 
%die Sensitivität $\vect{X}^{(\ell)}_{pp}$ (bzw. ihre Diagonale) liefert nur die
%Suchrichtung, nicht die Genauigkeit des Fixpunkts. Formal konvergiert die
%Schleife, wenn der Spektralradius der äußeren Iterationsabbildung $\rho < 1$
%bleibt, was nachfolgend numerisch berücksichtigt wird.


%\paragraph{Zusammenfassung der Konvergenzbedingungen: }
%\begin{center}
%	\begin{tabular}{lccc}
%		\toprule
%		Vorkonditionierung
%		& Linearisierter Spektralradius
%		& Skalierung mit $n_{pv}$
%		& Optimales $\omega$ \\
%		\midrule
%		Entkoppelt $\vect{D} = 2\diag(\vect{X}_{pp})$
%		& $\rho\bigl(\vect{I} - \omega\vect{D}^{-1}\vect{H}\bigr)$
%		& wächst
%		& $< 1$ (Dämpfung nötig) \\
%		Gekoppelt $\vect{A} = 2\vect{X}_{pp}$
%		& $|1-\omega|$
%		& konstant $=0$
%		& $\omega^\star = 1$ \\
%		\bottomrule
%	\end{tabular}
%\end{center}

\subsection{Algorithmus}
\label{subsec:methode_a_algorithmus}

\begin{algorithm}[H]
	\caption{Methode A: Äußere Q-Schleife mit innerem TPF (allgemeiner ZIP-Fall)}
	\label{alg:methode_a}
	\begin{algorithmic}[1]
		\REQUIRE $\vect{s}_{\text{PQ}}$, $P_k^{\spec}$ und $|V_k^{\spec}|$ für $k \in \Omega_{\text{PV}}$, $\Ydd$, $\Yds$, $\vect{v}_s$, ZIP-Koeffizienten $\alpha_P, \alpha_I, \alpha_Z$, Toleranzen $\varepsilon_v$, $\varepsilon_Q$, max.\ äußere Iterationen $L$
		\ENSURE $\vect{v}$, $Q_k$ für $k \in \Omega_{\text{PV}}$
		
		\STATE \textbf{Initialisierung:}
		\STATE $Q_k^{(0)} = 0$ für alle $k \in \Omega_{\text{PV}}$; \quad $\ell = 0$
		
		\STATE
		\STATE \textbf{Äußere Schleife:}
		\WHILE{$\max_k |\Delta|V_k|^2| \geq \varepsilon_Q$ \textbf{und} $\ell < L$}
		
		\STATE \textit{\% Leistungsvektor und Hilfsmatrizen aktualisieren}
		\STATE $s_k^{(\ell)} = P_k^{\spec} + \jj Q_k^{(\ell)}$ für alle $k \in \Omega_{\text{PV}}$
		\STATE $\vect{A}^{(\ell)} = \diag(\alpha_P \had \vect{s}^{*(\ell)})$
		\STATE $\vect{B}^{(\ell)} = \diag(\alpha_Z \had \vect{s}^{*(\ell)}) + \Ydd$
		\STATE $\vect{c}^{(\ell)} = \Yds\vect{v}_s + \alpha_I \had \vect{s}^{*(\ell)}$
		\STATE Berechne $(\vect{B}^{(\ell)})^{-1}$ \COMMENT{Nur wenn $\alpha_{Z,k} \neq 0$; sonst einmalig}
		\STATE $\vect{F}^{(\ell)} = -(\vect{B}^{(\ell)})^{-1}\vect{A}^{(\ell)}$; \quad $\vect{w}^{(\ell)} = -(\vect{B}^{(\ell)})^{-1}\vect{c}^{(\ell)}$
		\STATE $X_{kk}^{(\ell)} = \im\left([(\vect{B}^{(\ell)})^{-1}]_{kk}\right)$ für alle $k \in \Omega_{\text{PV}}$
		
		\STATE
		\STATE \textit{\% Innere Schleife: Standard-FPI mit konstantem $\vect{F}^{(\ell)}$, $\vect{w}^{(\ell)}$}
		\STATE $\vect{v}^{(0)} = \mathbf{1} + \jj\mathbf{0}$; \quad $n = 0$; \quad $\text{tol}_v = \infty$
		\WHILE{$\text{tol}_v \geq \varepsilon_v$ \textbf{und} $n < K$}
		\STATE $\vect{v}^{(n+1)} = \vect{F}^{(\ell)} \cdot \vect{v}_{(n)}^{*\eleminv} + \vect{w}^{(\ell)}$
		\STATE $\text{tol}_v = \max\|\vect{s}^{(\ell)} - \vect{s}^{(n+1)}_{\text{calc}}\|$; \quad $n = n + 1$
		\ENDWHILE
		\STATE $\vect{v}^{(\ell)} = \vect{v}^{(n)}$ \COMMENT{Konvergierte Lösung der inneren Schleife}
		
		\STATE
		\STATE \textit{\% Q-Korrektur}
		\FOR{jeden PV-Knoten $k \in \Omega_{\text{PV}}$}
		\STATE $\Delta|V_k|^2 = |V_k^{\spec}|^2 - |v_k^{(\ell)}|^2$
		\STATE $\Delta Q_k^{(\ell)} = \Delta|V_k|^2 \,/\, (2 \cdot X_{kk}^{(\ell)})$
		\STATE $Q_k^{(\ell+1)} = Q_k^{(\ell)} + \Delta Q_k^{(\ell)}$
		\STATE $Q_k^{(\ell+1)} = \max\big(Q_k^{\min},\, \min(Q_k^{\max},\, Q_k^{(\ell+1)})\big)$ \COMMENT{Q-Limits}
		\ENDFOR
		\STATE $\ell = \ell + 1$
		
		\ENDWHILE
		\RETURN $\vect{v}^{(\ell)}$, $\{Q_k^{(\ell)}\}_{k \in \Omega_{\text{PV}}}$
	\end{algorithmic}
\end{algorithm}

%\subsection{Rechentechnische Betrachtung}
%\label{subsec:methode_a_kosten}
%
%\paragraph{Fall $\alpha_{Z,k} = 0$ am PV-Knoten (häufiger Spezialfall).}
%Die Matrix $\vect{B}$ ist konstant und $\vect{B}^{-1}$ muss nur einmal berechnet werden. In der äußeren Schleife müssen lediglich die Diagonalmatrix $\vect{A}^{(\ell)}$, der Vektor $\vect{c}^{(\ell)}$ und die Produkte $\vect{F}^{(\ell)} = -\vect{B}^{-1}\vect{A}^{(\ell)}$, $\vect{w}^{(\ell)} = -\vect{B}^{-1}\vect{c}^{(\ell)}$ neu berechnet werden. Die Berechnung von $\vect{F}^{(\ell)}$ ist eine Matrix-Diagonalmatrix-Multiplikation (Skalierung der Spalten von $\vect{B}^{-1}$) mit Kosten $\mathcal{O}((b\phi)^2)$.
%
%\paragraph{Fall $\alpha_{Z,k} \neq 0$ (allgemeiner Fall).}
%Die Matrix $\vect{B}$ ändert sich, da ihr Diagonalelement am PV-Knoten $k$ den Wert $\alpha_{Z,k} s_k^{*(\ell)} + Y_{dd,kk}$ hat. Eine vollständige Neuberechnung von $\vect{B}^{-1}$ hat Kosten $\mathcal{O}((b\phi)^3)$. Stattdessen kann die Sherman-Morrison-Formel für Rang-1-Updates verwendet werden:
%\begin{equation}
%	(\vect{B}^{(\ell+1)})^{-1} = (\vect{B}^{(\ell)} + \delta_k \vect{e}_k \vect{e}_k^T)^{-1} = (\vect{B}^{(\ell)})^{-1} - \frac{(\vect{B}^{(\ell)})^{-1}\vect{e}_k \vect{e}_k^T (\vect{B}^{(\ell)})^{-1} \cdot \delta_k}{1 + \delta_k \vect{e}_k^T (\vect{B}^{(\ell)})^{-1} \vect{e}_k}
%	\label{eq:sherman_morrison}
%\end{equation}
%mit $\delta_k = \alpha_{Z,k}(s_k^{*(\ell+1)} - s_k^{*(\ell)}) = -\jj\,\alpha_{Z,k}\,\Delta Q_k^{(\ell)}$ und $\vect{e}_k$ dem $k$-ten Einheitsvektor. Die Kosten pro PV-Knoten reduzieren sich damit auf $\mathcal{O}((b\phi)^2)$.
%
%\paragraph{Konvergenz der inneren Schleife.}
%Da innerhalb der inneren Schleife $\vect{s}^{(\ell)}$ und damit $\vect{F}^{(\ell)}$ konstant sind, ist der Banach-Fixpunktsatz (\cref{prop:sam_kontraktion}) anwendbar. Die innere Schleife konvergiert garantiert, solange $\eta^{(\ell)} < 1$ mit:
%\begin{equation}
%	\eta^{(\ell)} = \left\|(\vect{B}^{(\ell)})^{-1} \diag(\alpha_P \had \vect{s}^{*(\ell)})\right\|_1 \cdot \frac{1}{v_{\min}^2}
%	\label{eq:eta_methode_a}
%\end{equation}

%%======================================================================
%\section{Methode B: Eingebettete Q-Korrektur}
%\label{sec:methode_b}
%%======================================================================
%
%\subsection{Grundidee}
%\label{subsec:methode_b_idee}
%
%Im Gegensatz zu Methode~A, die eine verschachtelte Schleifenstruktur erfordert, integriert Methode~B die PV-Knoten-Behandlung direkt in \textit{jede} FPI-Iteration. Die Grundidee besteht aus drei Schritten, die nach dem regulären FPI-Update ausgeführt werden:
%\begin{enumerate}
%	\item \textbf{Spannungsbetrags-Projektion:} Der an einem PV-Knoten berechnete Spannungszeiger wird in seiner Richtung (Winkel) beibehalten, aber im Betrag auf den Sollwert $|V_k^{\spec}|$ skaliert.
%	\item \textbf{Q-Rückrechnung:} Die Blindleistung wird aus der Netzgleichung (Kirchhoff'sches Knotengesetz) mit dem projizierten Spannungszeiger exakt zurückberechnet.
%	\item \textbf{Hilfsmatrizen-Update:} Die Matrizen $\vect{A}$, $\vect{B}$ (falls nötig), $\vect{c}$, $\vect{F}$, $\vect{w}$ werden mit dem neuen $Q_k$ aktualisiert.
%\end{enumerate}
%
%Dieser Ansatz erzwingt die Spannungsbedingung $|v_k| = |V_k^{\spec}|$ in jeder Iteration \textit{exakt}. Die Blindleistung ergibt sich direkt aus der Leistungsbilanz, ohne Linearisierung oder Sensitivitätsapproximation.
%
%\subsection{Mathematische Formulierung}
%\label{subsec:methode_b_formulierung}
%
%
%Die reguläre FPI-Iteration \eqref{eq:pv_fpi_basis} wird mit den aktuellen Matrizen $\vect{F}^{(n)}$, $\vect{w}^{(n)}$ durchgeführt:
%\begin{equation}
%	\vect{v}_{(n+1)}' = \vect{F}^{(n)} \cdot \vect{v}_{(n)}^{*\eleminv} + \vect{w}^{(n)}
%	\label{eq:methode_b_fpi}
%\end{equation}
%
%Der Strich in $\vect{v}_{(n+1)}'$ kennzeichnet, dass dies ein Zwischenergebnis ist, das an den PV-Knoten noch korrigiert wird. An PQ-Knoten gilt $v_{k,(n+1)} = v_{k,(n+1)}'$ ohne Korrektur.
%
%
%An jedem PV-Knoten $k \in \Omega_{\text{PV}}$ wird der Spannungszeiger auf den Sollbetrag projiziert:
%\begin{equation}
%	v_{k,(n+1)} = |V_k^{\spec}| \cdot \frac{v_{k,(n+1)}'}{|v_{k,(n+1)}'|}
%	\label{eq:methode_b_projektion}
%\end{equation}
%
%Dies ist eine reine Skalierungsoperation. Der resultierende Spannungszeiger hat den korrekten Betrag $|v_{k,(n+1)}| = |V_k^{\spec}|$ und einen aus der FPI-Iteration stammenden Winkel.
%
%
%Nach der Projektion ist die Spannung am PV-Knoten bekannt. Der zugehörige Strom ergibt sich aus der Admittanzmatrix:
%\begin{equation}
%	i_{k,(n+1)} = \sum_{m \in \Omega_d} Y_{dd,km} \cdot v_{m,(n+1)} + (\Yds)_k \cdot \vect{v}_s
%	\label{eq:methode_b_strom}
%\end{equation}
%
%Die komplexe Leistung am PV-Knoten ergibt sich zu:
%\begin{equation}
%	s_{k,(n+1)}^{\text{calc}} = v_{k,(n+1)} \cdot i_{k,(n+1)}^*
%	\label{eq:methode_b_leistung}
%\end{equation}
%
%Die Blindleistung ist der Imaginärteil:
%\begin{equation}
%	Q_k^{(n+1)} = \im\left(v_{k,(n+1)} \cdot i_{k,(n+1)}^*\right)
%	\label{eq:methode_b_Q}
%\end{equation}
%
%Die so berechnete Blindleistung die Netzgleichung \textit{exakt} erfüllt, im Gegensatz zur linearisierten Sensitivitätsformel \eqref{eq:methode_a_dQ} aus Methode~A.
%
%
%Der Leistungsvektor wird aktualisiert:
%\begin{equation}
%	s_k^{(n+1)} = P_k^{\spec} + \jj Q_k^{(n+1)}, \quad \forall k \in \Omega_{\text{PV}}
%	\label{eq:methode_b_s_update}
%\end{equation}
%
%Anschließend werden die Hilfsmatrizen für die nächste Iteration neu berechnet:
%\begin{align}
%	\vect{A}^{(n+1)} &= \diag(\alpha_P \had \vect{s}^{*(n+1)}) \label{eq:methode_b_A_update}\\
%	\vect{B}^{(n+1)} &= \diag(\alpha_Z \had \vect{s}^{*(n+1)}) + \Ydd \label{eq:methode_b_B_update}\\
%	\vect{c}^{(n+1)} &= \Yds\vect{v}_s + \alpha_I \had \vect{s}^{*(n+1)} \label{eq:methode_b_c_update}\\
%	\vect{F}^{(n+1)} &= -(\vect{B}^{(n+1)})^{-1}\vect{A}^{(n+1)} \label{eq:methode_b_F_update}\\
%	\vect{w}^{(n+1)} &= -(\vect{B}^{(n+1)})^{-1}\vect{c}^{(n+1)} \label{eq:methode_b_w_update}
%\end{align}
%
%
%In diesem häufigen Spezialfall sind $\vect{B}$, $\vect{B}^{-1}$ und $\vect{c}$ konstant. Es müssen lediglich die Einträge von $\vect{A}$ an den PV-Knoten-Positionen aktualisiert werden, und $\vect{F}$ ergibt sich durch Neuskalierung der entsprechenden Spalten von $\vect{B}^{-1}$. Die Berechnung von $\vect{w}$ entfällt, da $\vect{c}$ und $\vect{B}^{-1}$ unverändert bleiben.
%
%\subsection{Algorithmus}
%\label{subsec:methode_b_algorithmus}
%
%\begin{algorithm}[H]
%	\caption{Methode B: Eingebettete Q-Korrektur (allgemeiner ZIP-Fall)}
%	\label{alg:methode_b}
%	\begin{algorithmic}[1]
%		\REQUIRE $\vect{s}_{\text{PQ}}$, $P_k^{\spec}$ und $|V_k^{\spec}|$ für $k \in \Omega_{\text{PV}}$, $\Ydd$, $\Yds$, $\vect{v}_s$, ZIP-Koeffizienten, Toleranz $\varepsilon$, max.\ Iterationen $K$
%		\ENSURE $\vect{v}$, $Q_k$ für $k \in \Omega_{\text{PV}}$
%		
%		\STATE \textbf{Initialisierung:}
%		\STATE $Q_k^{(0)} = 0$ für alle $k \in \Omega_{\text{PV}}$
%		\STATE $\vect{v}^{(0)} = \mathbf{1} + \jj\mathbf{0}$; \quad $n = 0$; \quad $\text{tol} = \infty$
%		\STATE $\vect{s}^{(0)} = [\vect{s}_{\text{PQ}};\, P_k^{\spec} + \jj Q_k^{(0)}]$
%		\STATE Berechne $\vect{A}^{(0)}$, $\vect{B}^{(0)}$, $\vect{c}^{(0)}$, $(\vect{B}^{(0)})^{-1}$, $\vect{F}^{(0)}$, $\vect{w}^{(0)}$ gemäß \eqref{eq:pv_A}--\eqref{eq:pv_w}
%		
%		\STATE
%		\STATE \textbf{Hauptschleife:}
%		\WHILE{$\text{tol} \geq \varepsilon$ \textbf{und} $n < K$}
%		
%		\STATE \textit{\% Schritt 1: FPI-Schritt}
%		\STATE $\vect{v}_{(n+1)}' = \vect{F}^{(n)} \cdot \vect{v}_{(n)}^{*\eleminv} + \vect{w}^{(n)}$
%		
%		\STATE
%		\STATE \textit{\% Schritt 2: Spannungsbetrags-Projektion an PV-Knoten}
%		\FOR{jeden PV-Knoten $k \in \Omega_{\text{PV}}$}
%		\STATE $v_{k,(n+1)} = |V_k^{\spec}| \cdot v_{k,(n+1)}' \,/\, |v_{k,(n+1)}'|$
%		\ENDFOR
%		\STATE $v_{m,(n+1)} = v_{m,(n+1)}'$ für alle $m \notin \Omega_{\text{PV}}$
%		
%		\STATE
%		\STATE \textit{\% Schritt 3: Q-Rückrechnung aus Netzgleichung}
%		\FOR{jeden PV-Knoten $k \in \Omega_{\text{PV}}$}
%		\STATE $i_{k,(n+1)} = \sum_{m} Y_{dd,km} \cdot v_{m,(n+1)} + (\Yds)_k \cdot \vect{v}_s$
%		\STATE $Q_k^{(n+1)} = \im\left(v_{k,(n+1)} \cdot i_{k,(n+1)}^*\right)$
%		\STATE $Q_k^{(n+1)} = \max\big(Q_k^{\min},\, \min(Q_k^{\max},\, Q_k^{(n+1)})\big)$ \COMMENT{Q-Limits}
%		\ENDFOR
%		
%		\STATE
%		\STATE \textit{\% Schritt 4: Leistungsvektor und Hilfsmatrizen aktualisieren}
%		\STATE $s_k^{(n+1)} = P_k^{\spec} + \jj Q_k^{(n+1)}$ für alle $k \in \Omega_{\text{PV}}$
%		\STATE Aktualisiere $\vect{A}^{(n+1)}$ gemäß \eqref{eq:methode_b_A_update}
%		\IF{$\alpha_{Z,k} \neq 0$ für ein $k \in \Omega_{\text{PV}}$}
%		\STATE Aktualisiere $\vect{B}^{(n+1)}$ und $(\vect{B}^{(n+1)})^{-1}$ \COMMENT{Rang-1-Update}
%		\ENDIF
%		\IF{$\alpha_{I,k} \neq 0$ für ein $k \in \Omega_{\text{PV}}$}
%		\STATE Aktualisiere $\vect{c}^{(n+1)}$
%		\ENDIF
%		\STATE $\vect{F}^{(n+1)} = -(\vect{B}^{(n+1)})^{-1}\vect{A}^{(n+1)}$
%		\STATE $\vect{w}^{(n+1)} = -(\vect{B}^{(n+1)})^{-1}\vect{c}^{(n+1)}$
%		
%		\STATE
%		\STATE \textit{\% Konvergenz prüfen}
%		\STATE Berechne Leistungsmismatch an PQ-Knoten
%		\STATE $\text{tol} = \max\|\vect{s}^{\spec}_{\text{PQ}} - \vect{s}_{\text{PQ,calc}}^{(n+1)}\|$
%		\STATE $n = n + 1$
%		
%		\ENDWHILE
%		\RETURN $\vect{v}_{(n)}$, $\{Q_k^{(n)}\}_{k \in \Omega_{\text{PV}}}$
%	\end{algorithmic}
%\end{algorithm}
%
%
%\paragraph{Vergleich mit Methode A:}
%\begin{itemize}
%	\item Methode~B hat nur eine \textit{einzige Schleife}. Dies vereinfacht die Implementierung und vermeidet die vollständige Konvergenz einer inneren Schleife bei jedem äußeren Schritt.
%	\item Die Spannungsbedingung $|v_k| = |V_k^{\spec}|$ wird in \textit{jeder} Iteration exakt eingehalten.
%	\item Die Q-Bestimmung \eqref{eq:methode_b_Q} ist \textit{exakt} (keine Linearisierung), während Methode~A die Näherungsformel \eqref{eq:methode_a_dQ} verwendet.
%	\item Die Hilfsmatrizen müssen in jeder Iteration aktualisiert werden (im allgemeinen ZIP-Fall), während sie bei Methode~A innerhalb der inneren Schleife konstant bleiben.
%\end{itemize}




















%%======================================================================
%\section{Tensorformulierung der Methoden}
%\label{sec:tensor_formulierung}
%%======================================================================
%
%In diesem Abschnitt wird gezeigt, wie die in \cref{sec:methode_a,sec:methode_b} entwickelten Methoden für die gleichzeitige Berechnung von $\tau$ Lastflüssen im Tensor-Rahmen formuliert werden können.
%
%\subsection{Recap: Tensor-Reshape und allgemeine FPI}
%\label{subsec:tensor_recap}
%
%Gemäß \cref{sec:tpf_formulierungen} werden alle $\tau$ Leistungsvektoren als Spalten der Matrix $\dotS \in \C^{b\phi \times \tau}$ angeordnet. Die tensorisierte FPI in ihrer allgemeinen Form lautet:
%\begin{equation}
%	\dotV_{(n+1)} = \vect{F}^{(n)} \cdot \dotV_{(n)}^{*\eleminv} + \dot{W}^{(n)}
%	\label{eq:tensor_fpi_allgemein}
%\end{equation}
%
%wobei $\dot{W}^{(n)} = [\vect{w}^{(n)}, \ldots, \vect{w}^{(n)}] \in \C^{b\phi \times \tau}$ und $\dotV_{(n)}^{*\eleminv}$ die element-weise Anwendung der Operation ``konjugieren und Kehrwert bilden'' auf alle Elemente der Matrix darstellt.
%
%Im reinen PQ-Fall sind $\vect{F}$ und $\vect{w}$ konstant über alle Iterationen und alle $\tau$ Szenarien. Die Situation ändert sich bei PV-Knoten:
%
%\paragraph{Szenarioabhängigkeit.}
%Im Allgemeinen haben verschiedene Szenarien $i$ und $j$ unterschiedliche Lastbedingungen und damit unterschiedliche Q-Werte am selben PV-Knoten:
%\begin{equation}
%	Q_k^{(n)}[i] \neq Q_k^{(n)}[j] \quad \Rightarrow \quad s_k^{*(n)}[i] \neq s_k^{*(n)}[j]
%	\label{eq:szenarioabh}
%\end{equation}
%
%Dies bedeutet, dass der Leistungsvektor für verschiedene Szenarien unterschiedlich ist. Folglich wären im allgemeinen Fall die Matrizen $\vect{A}$, $\vect{B}$, $\vect{F}$, $\vect{w}$ für \textit{jedes Szenario} unterschiedlich -- was die Tensor-Parallelisierung zunichte machen würde.
%
%\paragraph{Die rettende Eigenschaft: Nur $\vect{A}$ ist szenarioabhängig (für $\alpha_{Z,k} = 0$).}
%Wenn am PV-Knoten $\alpha_{Z,k} = 0$ gilt, dann hängt $\vect{B}$ nicht von $\vect{s}$ ab und ist daher für alle Szenarien und Iterationen \textit{identisch}. Ebenso ist $\vect{c}$ konstant falls $\alpha_{I,k} = 0$. Der Tensor-Trick bleibt erhalten, indem die szenarioabhängige Matrix $\vect{A}^{(n)}$ in die Hadamard-Struktur aufgelöst wird.
%
%\subsection{Umformulierung für den Tensor-Fall}
%\label{subsec:tensor_umformulierung}
%
%Für den Spezialfall $\alpha_{Z,k} = \alpha_{I,k} = 0$ am PV-Knoten schreibt sich die Iteration \eqref{eq:pv_fpi_explizit2} als:
%\begin{equation}
%	\vect{v}_{(n+1)} = -\vect{B}^{-1} \cdot \left[(\alpha_P \had \vect{s}^{*(n)}) \had \vect{v}_{(n)}^{*\eleminv}\right] - \vect{B}^{-1}\vect{c}
%	\label{eq:tensor_umform1}
%\end{equation}
%
%wobei $\vect{B}^{-1}$ und $\vect{c}$ konstant sind. Im Tensor-Format:
%\begin{equation}
%	\dotV_{(n+1)} = -\vect{B}^{-1} \cdot \underbrace{\left[(\alpha_P \had \dotS^{*(n)}) \had \dotV_{(n)}^{*\eleminv}\right]}_{\in \C^{b\phi \times \tau}} + \dot{W}
%	\label{eq:tensor_umform2}
%\end{equation}
%
%Die Matrix $\dotS^{*(n)} \in \C^{b\phi \times \tau}$ enthält die szenario- und iterationsabhängigen Leistungswerte. Entscheidend ist: Die \textit{linke} Matrix $\vect{B}^{-1}$ ist für alle $\tau$ Szenarien und alle Iterationen \textit{identisch}. Die Multiplikation $\vect{B}^{-1} \cdot (\ldots)$ bleibt eine einzige Matrix-Matrix-Multiplikation -- unabhängig davon, ob sich $\dotS^{(n)}$ zwischen den Iterationen ändert.
%
%\paragraph{Allgemeiner ZIP-Fall ($\alpha_{Z,k} \neq 0$).}
%Wenn $\vect{B}$ szenarioabhängig wird, ist die Tensor-Parallelisierung in der dense-Formulierung nicht mehr direkt möglich, da jedes Szenario seine eigene Matrix $(\vect{B}^{(n)}[i])^{-1}$ benötigen würde. In diesem Fall muss auf die sparse-Formulierung zurückgegriffen werden, bei der die Block-Diagonalmatrix $\dotM$ (vgl.\ \eqref{eq:block_diagonal}) ohnehin szenarioabhängige Blöcke enthält. Die einmalige Faktorisierung geht dann zwar verloren, aber die Block-Diagonalstruktur ermöglicht weiterhin eine effiziente parallele Lösung.
%
%In der Praxis ist dieser Fall selten, da Generatoren als leistungskonstant modelliert werden. Für den Rest dieses Abschnitts wird daher $\alpha_{Z,k} = \alpha_{I,k} = 0$ am PV-Knoten vorausgesetzt.

%\subsection{Tensorformulierung}
%\label{subsec:tensor_methode_a}
%
%Bei Methode~A werden alle $\tau$ Lastflüsse \textit{gleichzeitig} in der inneren Schleife berechnet. Die Q-Korrektur in der äußeren Schleife wird element-weise über alle $\tau$ Szenarien ausgeführt:
%
%\paragraph{Innere Schleife (Matrix-Matrix-Multiplikation)}
%In jeder äußeren Iteration $\ell$ wird $\vect{F}^{(\ell)}$ aktualisiert und die innere Schleife parallelisiert:
%\begin{equation}
%	\dotV_{(n+1)} = \vect{F}^{(\ell)} \cdot \dotV_{(n)}^{*\eleminv} + \dot{W}
%	\label{eq:tensor_a_inner}
%\end{equation}
%
%Dies ist eine einzige Multiplikation $\vect{F}^{(\ell)} \in \C^{b\phi \times b\phi}$ mit $\dotV_{(n)}^{*\eleminv} \in \C^{b\phi \times \tau}$ für alle $\tau$ Lastflüsse gleichzeitig. Da $\vect{F}^{(\ell)}$ innerhalb der inneren Schleife konstant ist, ist dies exakt die TPF-Operation aus \cref{alg:tpf_dense}.
%
%%\paragraph{Bemerkung zur Szenarioabhängigkeit.}
%Im Tensor-Fall hat jedes Szenario $i$ einen anderen Leistungsvektor $\dotS[:,i]$ und damit nach der äußeren Q-Korrektur einen anderen Q-Wert $Q_k^{(\ell)}[i]$. Die aktualisierte Matrix $\vect{A}^{(\ell)}$ müsste eigentlich szenarioabhängig sein. Für Methode~A im Tensor-Setting wird daher die innere Schleife umformuliert als:
%\begin{equation}
%	\dotV_{(n+1)} = -\vect{B}^{-1} \cdot \left[(\alpha_P \had \dotS^{*(\ell)}) \had \dotV_{(n)}^{*\eleminv}\right] + \dot{W}
%	\label{eq:tensor_a_inner_explizit}
%\end{equation}
%
%wobei $\dotS^{*(\ell)} \in \C^{b\phi \times \tau}$ für jedes Szenario den aktuellen Leistungsvektor enthält und innerhalb der inneren Schleife \textit{fest} ist. Die Operation $(\alpha_P \had \dotS^{*(\ell)}) \had \dotV_{(n)}^{*\eleminv}$ ist eine rein element-weise Operation auf der $b\phi \times \tau$-Matrix und die anschließende Multiplikation mit $\vect{B}^{-1}$ ist eine Matrix-Matrix-Multiplikation. Der Tensor-Vorteil bleibt vollständig erhalten.
%
%Die Äußere Schleife ich über einfache elementweise rechenoperationen durch
%\begin{equation}
%	\Delta Q[k,:] = \frac{|V_k^{\spec}|^2 - |\dotV^{(\ell)}_{\conv}[k,:]|^2}{2 \cdot X_{kk}}, \quad \forall k \in \Omega_{\text{PV}}
%	\label{eq:tensor_a_dQ}
%\end{equation}
%tensorisierbar.
%Die Division durch den Skalar $2X_{kk}$ wirkt gleichzeitig auf alle $\tau$ Szenarien.
%
%\subsection{Tensorformulierung von Methode B}
%\label{subsec:tensor_methode_b}
%
%Methode~B erfordert keine verschachtelte Schleifenstruktur und ist besonders gut für die Tensorisierung geeignet. Die vier Schritte (\cref{alg:methode_b}) werden direkt auf den Tensor-Fall übertragen:
%
%\paragraph{Schritt 1: FPI (Matrix-Matrix-Multiplikation).}
%\begin{equation}
%	\dotV_{(n+1)}' = -\vect{B}^{-1} \cdot \left[(\alpha_P \had \dotS^{*(n)}) \had \dotV_{(n)}^{*\eleminv}\right] + \dot{W}
%	\label{eq:tensor_b_fpi}
%\end{equation}
%
%Die Matrix $\vect{B}^{-1}$ bleibt über alle Iterationen und Szenarien konstant (für $\alpha_{Z,k} = 0$). Die szenario- und iterationsabhängige Information steckt vollständig in $\dotS^{*(n)}$, welches sich an den PV-Knoten-Zeilen zwischen den Iterationen ändert.
%
%\paragraph{Schritt 2: Projektion (element-weise).}
%\begin{equation}
%	\dotV_{(n+1)}[k,:] = |V_k^{\spec}| \cdot \frac{\dotV_{(n+1)}'[k,:]}{|\dotV_{(n+1)}'[k,:]|}, \quad \forall k \in \Omega_{\text{PV}}
%	\label{eq:tensor_b_proj}
%\end{equation}
%
%Betragsberechnung und Division wirken element-weise auf die Zeile $[k,:]$ und damit gleichzeitig auf alle $\tau$ Szenarien.
%
%\paragraph{Schritt 3: Stromberechnung und Q-Rückrechnung.}
%\begin{equation}
%	\dot{I}_{\text{PV}} = \Ydd[\Omega_{\text{PV}}, :] \cdot \dotV_{(n+1)} + \Yds[\Omega_{\text{PV}}, :] \cdot \vect{v}_s
%	\label{eq:tensor_b_strom_matrix}
%\end{equation}
%
%Dies ist eine Matrix-Matrix-Multiplikation der Dimension $(n_{\text{PV}} \times b\phi) \cdot (b\phi \times \tau) = (n_{\text{PV}} \times \tau)$. Die Blindleistung ergibt sich element-weise:
%\begin{equation}
%	\dot{Q}_{\text{PV}}^{(n+1)} = \im\left(\dotV_{(n+1)}[\Omega_{\text{PV}}, :] \had \dot{I}_{\text{PV}}^*\right) \in \R^{n_{\text{PV}} \times \tau}
%	\label{eq:tensor_b_Q_matrix}
%\end{equation}
%
%\paragraph{Schritt 4: $\dotS$-Update (element-weise).}
%\begin{equation}
%	\dotS^{(n+1)}[k,:] = P_k^{\spec} + \jj \dot{Q}_{\text{PV}}^{(n+1)}[k,:], \quad \forall k \in \Omega_{\text{PV}}
%	\label{eq:tensor_b_s_update}
%\end{equation}
%
%\subsection{Gesamtvergleich der Tensor-Operationen}
%\label{subsec:tensor_vergleich}
%
%\begin{table}[H]
%	\centering
%	\caption{Rechentechnischer Vergleich der Tensor-Operationen pro Iteration (Spezialfall $\alpha_{Z,k} = \alpha_{I,k} = 0$ am PV-Knoten).}
%	\label{tab:tensor_vergleich}
%	\begin{tabular}{@{}p{4cm}p{4.5cm}cc@{}}
%		\toprule
%		\textbf{Operation} & \textbf{Dimension} & \textbf{Kosten} & \textbf{Gleich wie PQ?} \\
%		\midrule
%		$\vect{B}^{-1} \cdot (\ldots)$ & $(b\phi)^2 \cdot \tau$ & Dominant & Ja \\
%		Projektion & $n_{\text{PV}} \cdot \tau$ (element-weise) & Minimal & -- \\
%		$Y_{dd, \text{PV}} \cdot \dotV$ & $n_{\text{PV}} \cdot b\phi \cdot \tau$ & Gering$^*$ & -- \\
%		Q-Berechnung & $n_{\text{PV}} \cdot \tau$ (element-weise) & Minimal & -- \\
%		$\dotS$-Update & $n_{\text{PV}} \cdot \tau$ (element-weise) & Minimal & -- \\
%		\midrule
%		$\vect{B}^{-1}$ \textbf{konstant?} & & \textbf{Ja} & \\
%		\textbf{Tensor-Parallelisierung?} & & \textbf{Ja} & \\
%		\bottomrule
%	\end{tabular}\\[3pt]
%	{\footnotesize $^*$ $n_{\text{PV}} \ll b\phi$ in typischen Verteilnetzen, daher vernachlässigbar gegenüber der Hauptoperation.}
%\end{table}
%
%\subsection{Zusammenfassung: Warum die Parallelisierung erhalten bleibt}
%\label{subsec:tensor_zusammenfassung}
%
%Die Analyse zeigt, dass auch bei Verwendung der allgemeinen FPI-Formulierung \eqref{eq:pv_fpi_basis} der Tensor-Parallelisierungsvorteil erhalten bleibt, sofern die folgende Bedingung erfüllt ist:
%
%\begin{equation}
%	\alpha_{Z,k} = 0, \quad \forall k \in \Omega_{\text{PV}}
%	\label{eq:tensor_bedingung}
%\end{equation}
%
%Unter dieser Bedingung -- die für leistungskonstante Generatoren stets erfüllt ist -- gilt:
%\begin{enumerate}
%	\item Die Matrix $\vect{B}^{-1}$ ist \textbf{konstant über alle Iterationen und alle $\tau$ Szenarien}. Ihre einmalige Berechnung (oder Faktorisierung im sparse-Fall) wird über die gesamte Rechnung wiederverwendet.
%	\item Die teure Operation (Matrix-Matrix-Multiplikation $\vect{B}^{-1} \cdot (\ldots)$) ist \textbf{strukturell identisch} zum reinen PQ-Fall und profitiert uneingeschränkt von BLAS/GPU-Parallelisierung.
%	\item Die Änderung des Leistungsvektors $\dotS^{(n)}$ an PV-Knoten betrifft nur die \textbf{rechte Seite} des Produkts (den Vektor, der mit $\vect{B}^{-1}$ multipliziert wird) und nicht die Matrix selbst.
%	\item Die zusätzlichen Operationen (Projektion, Stromberechnung, Q-Update) sind entweder element-weise oder von deutlich geringerer Dimension ($n_{\text{PV}} \ll b\phi$) und damit vernachlässigbar.
%	\item Die einzige Einschränkung ist eine erhöhte Iterationszahl: Methode~B benötigt typischerweise 8--15 statt 5--10 Iterationen, Methode~A benötigt 15--50 FPI-Evaluierungen. Der Speedup pro Iteration bleibt jedoch unverändert.
%\end{enumerate}
%
%Wenn $\alpha_{Z,k} \neq 0$ am PV-Knoten gilt, geht der Tensor-Vorteil in der dense-Formulierung teilweise verloren, da $\vect{B}^{-1}$ szenarioabhängig wird. In diesem Fall bietet die sparse-Formulierung (\cref{subsec:tpf_sparse}) mit szenarioabhängigen Diagonalblöcken in $\dotM$ eine Alternative, wobei die Block-Diagonalstruktur eine parallele Lösung ermöglicht.
%
%Die quantitative Bewertung des resultierenden Gesamtspeedups gegenüber konventionellen Verfahren ist Gegenstand von \cref{ch:ergebnisse}.




%
%\subsection{Herleitung der Q-Korrekturformel}
%\label{subsec:methode_a_herleitung}
%
%Für die Q-Korrektur in Schritt~5 wird eine Beziehung zwischen $\Delta Q_k$ und dem
%Spannungsmismatch $\Delta|V_k|^2$ benötigt. Sie folgt unmittelbar aus der bereits
%aufgestellten Netzgleichung, ohne dass ein zusätzliches Ersatzschaltbild postuliert
%werden muss.
%
%%----------------------------------------------------------------------
%Ausgangspunkt ist die Fixpunktform \eqref{eq:fpi_dense}, die sich mit der
%Bus-Impedanzmatrix $\vect{Z}^{(\ell)}_B = (\vect{B}^{(\ell)})^{-1}$ als
%\begin{equation}
%	\vect{v} \;=\; \vect{Z}^{(\ell)}_B\,\vect{i} \;+\; \vect{w},
%	\qquad
%	i_k \;=\; \frac{s_k^*}{v_k^*},
%	\label{eq:netzgleichung_zb}
%\end{equation}
%schreiben lässt. Diese Beziehung ist in den Strominjektionen $\vect{i}$ exakt und
%linear; die gesamte Nichtlinearität des Lastflussproblems steckt allein in der
%Abhängigkeit $i_k(v_k)$. Spaltet man Zeile~$k$ in den Eigen- und den Fremdanteil auf,
%\begin{equation}
%	v_k \;=\; \underbrace{w_k + \sum_{j\neq k} Z^{(\ell)}_{B,kj}\,i_j}_{E_{th,k}}
%	\;+\; \underbrace{Z^{(\ell)}_{B,kk}}_{Z_{th,k}}\;i_k,
%	\label{eq:thevenin_identifikation}
%\end{equation}
%so ist das Thévenin-Äquivalent am Knoten~$k$ unmittelbar identifiziert: Die
%Leerlaufspannung $E_{th,k}$ ist die Überlagerung von Slack-Einspeisung und allen
%übrigen Knotenströmen, die Thévenin-Impedanz ist das Diagonalelement
%\begin{equation}
%	Z^{(\ell)}_{B,kk} \;=\; \bigl[(\vect{B}^{(\ell)})^{-1}\bigr]_{kk}
%	\;=\; R^{(\ell)}_{kk} + \jj\,X^{(\ell)}_{kk}.
%	\label{eq:diagonalelement}
%\end{equation}
%Ein Ersatzschaltbild muss somit nicht als Modellannahme eingeführt werden; es ist
%eine Umschreibung von~\eqref{eq:netzgleichung_zb}. Insbesondere braucht $E_{th,k}$
%nie explizit berechnet zu werden, da es in der folgenden Differenzenbildung entfällt.
%
%%----------------------------------------------------------------------
%Da $\vect{Z}^{(\ell)}_B$ und $\vect{w}$ im konstanten Lastfall nach
%\cref{subsec:konsequenz_analyse} unabhängig von $\vect{Q}$ sind, folgt aus
%\eqref{eq:netzgleichung_zb} für eine Störung der Strominjektionen exakt
%\begin{equation}
%	\Delta\vect{v} \;=\; \vect{Z}^{(\ell)}_B\,\Delta\vect{i}.
%	\label{eq:delta_v_exakt}
%\end{equation}
%Zwei Näherungen führen von hier zur Korrekturformel. Erstens wird die
%Strominjektion am PV-Knoten um den aktuellen Arbeitspunkt $v^{(\ell)}_k$
%linearisiert; mit $s_k = P_k + \jj Q_k$ in Erzeugerkonvention und festgehaltenem
%$P_k$ gilt
%\begin{equation}
%	\Delta i_k \;\approx\; \frac{-\jj\,\Delta Q_k}{(v^{(\ell)}_k)^*}.
%	\label{eq:strom_linear}
%\end{equation}
%Zweitens wird die Rückwirkung auf die übrigen Knoten vernachlässigt
%($\Delta i_j = 0$ für $j \neq k$), obwohl sich bei konstanter Leistung auch die
%Lastströme mit den Spannungen ändern. Beide Näherungen betreffen ausschließlich die
%Schrittweite der äußeren Korrektur, nicht deren Fixpunkt: Das Abbruchkriterium
%bewertet die Spannungsbeträge der exakten TPF-Lösung, sodass ein Fehler in der
%Sensitivität lediglich die Konvergenzgeschwindigkeit beeinflusst.
%
%Einsetzen von~\eqref{eq:strom_linear} in~\eqref{eq:delta_v_exakt} liefert
%\begin{equation}
%	\Delta v_k \;=\; -\jj\,Z^{(\ell)}_{B,kk}\,\frac{\Delta Q_k}{(v^{(\ell)}_k)^*}.
%	\label{eq:spannung_skalar}
%\end{equation}
%Mit $-\jj\,(R^{(\ell)}_{kk} + \jj X^{(\ell)}_{kk}) = X^{(\ell)}_{kk} - \jj R^{(\ell)}_{kk}$
%und $(v^{(\ell)}_k)^{*-1} = e^{\jj\theta_k}/|v^{(\ell)}_k|$ ergibt die
%Ausmultiplikation
%\begin{equation}
%	\Delta v_k \;=\; \frac{\Delta Q_k}{|v^{(\ell)}_k|}
%	\bigl(X^{(\ell)}_{kk} - \jj R^{(\ell)}_{kk}\bigr)\,e^{\jj\theta_k}.
%	\label{eq:spannung_ausmult}
%\end{equation}